\documentclass[
  aps,
  prb,
  preprint,
  superscriptaddress,
  longbibliography,
  nofootinbib
]{revtex4-2}

\usepackage[T1]{fontenc}
\usepackage[utf8]{inputenc}
\usepackage{amsmath,amssymb,bm}
\usepackage{dsfont}
\usepackage{mathtools}
\usepackage{booktabs}
\usepackage{graphicx}
\usepackage{microtype}
\usepackage{xcolor}
\usepackage{xr-hyper}
\externaldocument[M-]{paper}
\usepackage[colorlinks=true,allcolors=blue!45!black]{hyperref}

\newcommand{\dd}{\mathrm{d}}
\newcommand{\ee}{\mathrm{e}}
\newcommand{\ii}{\mathrm{i}}
\newcommand{\Tr}{\operatorname{Tr}}
\newcommand{\avg}[1]{\left\langle #1 \right\rangle}
\newcommand{\bR}{\mathbf{R}}
\newcommand{\br}{\mathbf{r}}
\newcommand{\bk}{\mathbf{k}}
\newcommand{\bp}{\mathbf{p}}
\newcommand{\bv}{\mathbf{v}}
\newcommand{\bnabla}{\boldsymbol{\nabla}}
\newcommand{\vF}{v_{\mathrm{F}}}
\newcommand{\vg}{v_{\mathrm{g}}}
\newcommand{\tautr}{\tau_{\mathrm{tr}}}
\newcommand{\elltr}{\ell_{\mathrm{tr}}}
\newcommand{\DOS}{N_1}
\newcommand{\DN}{D_{\mathrm{N}}}
\newcommand{\DE}{D(E)}
\newcommand{\LD}{L_D}
\newcommand{\Icoll}{I_{\mathrm{coll}}}
\newcommand{\fL}{f_L}
\newcommand{\fT}{f_T}
\newcommand{\qp}{\mathrm{qp}}
% --- merged in for the unified document (derivation-only macros + cleveref) ---
\usepackage[capitalize]{cleveref}
\crefname{equation}{Eq.}{Eqs.}
\Crefname{equation}{Eq.}{Eqs.}
\crefname{table}{Table}{Tables}
\Crefname{table}{Table}{Tables}
\crefname{section}{Sec.}{Secs.}
\Crefname{section}{Sec.}{Secs.}
\crefname{subsection}{Sec.}{Secs.}
\Crefname{subsection}{Sec.}{Secs.}
\crefname{subsubsection}{Sec.}{Secs.}
\Crefname{subsubsection}{Sec.}{Secs.}
\newcommand{\bA}{\mathbf{A}}
\newcommand{\hatg}{\hat{g}}
\newcommand{\checkg}{\check{g}}
\newcommand{\hatD}{\hat{\Delta}}
\providecommand{\openone}{}\renewcommand{\openone}{\mathds{1}}
\newcommand{\Real}{\operatorname{Re}}
\newcommand{\Imag}{\operatorname{Im}}

\renewcommand{\thesection}{S\Roman{section}}
\renewcommand{\theequation}{S\arabic{equation}}

\begin{document}

\title{Supplemental Material for ``Which Diffusion Equation for Nonequilibrium Superconducting Quasiparticles?''}

\author{Soren Ormseth}
\affiliation{Department of Physics, University of Wisconsin--Madison, Madison, Wisconsin 53706, USA}

\date{\today}

\maketitle

\section{Detailed dirty-limit derivation}\label{app:derivation}


\subsection*{Statement of the result}

Starting from the matrix Keldysh--Eilenberger equation, derive a scalar
Boltzmann-like equation for the positive-energy quasiparticle occupation
\(f(E,\br,t)\) valid in the dirty limit \(\tau\Delta\ll 1\). Target form:
\begin{equation}
  \DOS(E,\br)\,\partial_t f
  + \DOS(E,\br)\,\frac{\Delta}{E}\dot{\Delta}\,\partial_E f
  =
  \bnabla\!\cdot\!\bigl[\DN\,\mathcal D_L(E,\br)\,\bnabla f\bigr]
  + \Icoll^{\mathrm{inel}}[f],
  \label{eq:target}
\end{equation}
with \(\mathcal D_L=1\) above the local gap edge and \(0\) below it.
This is the \((p,q)=(1,0)\) operator of the main text. The elastic impurity
collision integral does not appear on the right-hand side: in the dirty limit it
gets used up to slave the first angular harmonic to the gradient of the zeroth,
generating the diffusion operator on the right of \cref{eq:target}. The only
collision integral that survives explicitly is inelastic
(electron-phonon, photon, etc.).

The whole calculation is organized as a sequence of well-defined reductions:
\begin{enumerate}
  \item Matrix Keldysh--Eilenberger \(\to\) matrix Keldysh--Usadel
        (angular reduction at the matrix level).
  \item Matrix Keldysh--Usadel \(\to\) coupled \((\fL,\fT)\) kinetic equations
        (Keldysh projection with the Schmid--Sch\"on--Larkin--Ovchinnikov
        distribution-matrix ansatz).
  \item Coupled \((\fL,\fT)\) \(\to\) scalar \(f(E,\br,t)\) equation
        (\(\fT\!\to\!0\) reduction, change of variable).
\end{enumerate}
Each step has a clear cost: which physics is dropped, and which symmetries are
used. The setup below pins everything down before we start computing.

\subsection{Conventions}
\label{sec:conventions}

We collect here the conventions used throughout the main text and this
appendix.

\paragraph{Nambu space.} \(\tau_1, \tau_2, \tau_3\) are Pauli matrices in the
particle-hole (Nambu) basis. We use
\[
  \tau_3=\mathrm{diag}(+1,-1),
  \qquad
  \hatD(\br,t)=-\ii\,\Delta(\br,t)\,\tau_2 ,
\]
for the particle-hole-symmetric singlet gap in the real-axis convention. The
anomalous component is placed along \(\tau_2\) so that the gap matrix matches
the off-diagonal structure of the BCS quasiclassical propagators
\cref{eq:gRA_BCS}; with this choice the coherent spectral operator
\(\hat L_0=E\tau_3-\hatD\) of \cref{sec:K_block} takes the explicit form
\[
  \hat L_0=E\tau_3+\ii\,\Delta\,\tau_2 ,
\]
which is proportional to the retarded propagator \(\hatg^{R}\) on the BCS
branch, so that \([\hat L_0,\hatg^{R}]=0\) and \(\hatg^{R}\) is the static
spectral solution (verified explicitly in \cref{sec:moyal_spectral}). The
alternative \(\hatD=\Delta\tau_1\) does not commute with the
\(\tau_2\)-anomalous propagators \cref{eq:gRA_BCS} and is not used here.

\paragraph{Quasiclassical Green's functions.} Energy-integrated propagators on
the Fermi surface,
\[
  \hatg^{R/A/K}(E,\hat\bp,\br,t)
  =
  \frac{\ii}{\pi}\!\int\!\dd\xi_\bp\,
  \tau_3\,\mathcal{G}^{R/A/K}\!\bigl(\bp,\br,t\bigr),
\]
are \(2\times 2\) matrices in Nambu space depending on energy \(E\), momentum
direction \(\hat\bp\) (unit vector on the Fermi surface), spatial
coordinate \(\br\) and time \(t\). The \(\tau_3\) fold in the definition is
the standard quasiclassical convention: it places the BCS density of states
on the \(\tau_3\) component of \(\hatg^{R}\) rather than on the identity,
and it converts the Nambu retarded--advanced symmetry
\(\hat{\mathcal G}^{A}=[\hat{\mathcal G}^{R}]^{\dagger}\) into the
conjugation relation \(\hatg^{A}=-\tau_3\hatg^{R\dagger}\tau_3\) of the
main text (\cref{M-eq:bcs_advanced_matrix}). The retarded/advanced/Keldysh blocks satisfy
the Keldysh quasiclassical normalization
\begin{equation}
  \hatg^{R}\circ\hatg^{R}=\hatg^{A}\circ\hatg^{A}=\openone,
  \qquad
  \hatg^{R}\circ\hatg^{K}+\hatg^{K}\circ\hatg^{A}=0,
  \label{eq:norm_RAK}
\end{equation}
where ``\(\circ\)'' is the Wigner-representation Moyal product
\(A\circ B=A\,\ee^{(\ii/2)(\overleftarrow{\partial}\!\cdot\overrightarrow{\partial})}B\)
in the relative-coordinate Fourier conjugate variables. To leading gradient
order it reduces to the ordinary product. We will work to that order
throughout (the Poisson-bracket corrections produce the time-derivative and
spectral-flow terms explicitly in \cref{sec:scalar_reduction}).

\paragraph{Distribution matrix.} The Keldysh component is parameterized as
\begin{equation}
  \hatg_{0}^{K}
  =
  \hatg_{0}^{R}\circ \hat h - \hat h\circ \hatg_{0}^{A},
  \qquad
  \hat h(E,\br,t)=\fL(E,\br,t)\,\openone+\fT(E,\br,t)\,\tau_3 ,
  \label{eq:dist_matrix_ansatz}
\end{equation}
with subscript ``\(0\)'' denoting the Fermi-surface average (introduced
properly in \cref{sec:angular_ansatz}). In equilibrium
\(\fL^{(0)}(E)=\tanh(E/2k_{\mathrm B}T)\), \(\fT^{(0)}=0\), and the positive-energy
quasiparticle occupation is
\(f(E,\br,t)=(1-\fL(E,\br,t))/2\), with
\(f^{(0)}(E)=(\ee^{E/k_{\mathrm B}T}+1)^{-1}\). The two-component reduction of the
generalized distribution matrix to the form
\cref{eq:dist_matrix_ansatz} is the
Schmid--Sch\"on--Larkin--Ovchinnikov form
(Schmid and Sch\"on~\cite{schmidschon1975}; Schmid~\cite{schmid1980}; Larkin and Ovchinnikov~\cite{larkin1986};
Kopnin~\cite{kopnin2001}, \S 9.2).

\paragraph{Spectral functions.} For a real, time-dependent gap on the
positive-energy quasiparticle branch,
\begin{align}
  \DOS(E,\br) &= \Real\,\frac{E}{\sqrt{(E+\ii 0^+)^2-\Delta^2(\br)}}
              = \frac{E}{\sqrt{E^2-\Delta^2(\br)}}\,\theta(E-\Delta), \\
  N_2(E,\br) &= \Imag\,\frac{-\ii\Delta(\br)}{\sqrt{(E+\ii 0^+)^2-\Delta^2}}
              = \frac{\Delta(\br)}{\sqrt{E^2-\Delta^2(\br)}}\,\theta(E-\Delta),
\end{align}
satisfying \(\DOS^2-N_2^2=1\) on the quasiparticle branch. The retarded
component decomposes in our basis as
\(\hatg^{R}=g^{R}\tau_3+\ii f^{R}\tau_2\) with
\(g^{R}=E/\sqrt{(E+\ii 0^+)^2-\Delta^2}\) and
\(f^{R}=\Delta/\sqrt{(E+\ii 0^+)^2-\Delta^2}\); above the gap and in the
\(\eta\to 0^+\) limit, \(\Real g^{R}=\DOS\), \(\Real f^{R}=N_2\), and
\(g^{A}=(g^{R})^*\), \(f^{A}=(f^{R})^*\).

\paragraph{Diffusion constant.} \(\DN=\vF^2\tau/d\) is the normal-state
diffusion constant in \(d\) spatial dimensions; we will write \(d=3\) when a
specific dimension is needed but otherwise carry \(d\) symbolically.
\(\tau\) is the elastic impurity transport time entering the \(s\)-wave Born
self-energy below.

\subsection{Starting equation: matrix Keldysh--Eilenberger}
\label{sec:keldysh_eilenberger}

The matrix quasiclassical equation of motion in the Keldysh representation
is~\cite{kopnin2001,larkin1986,rammer1986,belzig1999}
\begin{equation}
  \boxed{\,
  \ii\hbar\vF\,\hat\bp\!\cdot\!\widetilde{\bnabla}_{\!\br}\,\checkg
  +
  \bigl[\,
    E\,\tau_3 - \hatD(\br,t) - \check{\Sigma}(\hat\bp,\br,t;E),
    \;\checkg
  \,\bigr]_{\circ}
  =
  0,
  \qquad
  \checkg\circ\checkg=\openone,
  \,}
  \label{eq:keldysh_eilenberger}
\end{equation}
The factor of \(\ii\) on the streaming term is the standard Belzig--Kopnin
real-frequency convention; with it, the impurity self-energy below is purely
imaginary and the algebraic identities used in the slaving step
(\cref{sec:slaving}) produce a real coefficient. Other references absorb the
\(\ii\) into the streaming and put it on the impurity self-energy instead;
that is a global rescaling of the equation, not a different physical theory,
and the matrix Keldysh--Usadel form
\cref{eq:matrix_keldysh_usadel} is invariant.

Here \(\checkg\) is the upper-triangular \(2\times 2\) matrix in Keldysh space
\[
  \checkg
  =
  \begin{pmatrix}
    \hatg^{R} & \hatg^{K} \\
    0 & \hatg^{A}
  \end{pmatrix},
  \qquad
  \check\Sigma
  =
  \begin{pmatrix}
    \hat\Sigma^{R} & \hat\Sigma^{K} \\
    0 & \hat\Sigma^{A}
  \end{pmatrix},
\]
and \(\widetilde{\bnabla}_{\!\br}X\equiv\bnabla_{\!\br}X
-(\ii e/\hbar c)[\tau_3\bA(\br,t),X]\) is the gauge-covariant derivative.
The square-bracket commutator is taken with the Moyal product, which to
leading gradient order reduces to the ordinary commutator plus the Poisson-bracket
correction
\([A,B]_{\circ}\to[A,B]+(\ii\hbar/2)\{A,B\}_{\mathrm{P.B.}}\)
in mixed (Wigner) variables. The Poisson bracket on the
\(E\tau_3\) and \(\hatD\) pieces is the candidate source of the
time-derivative \(\partial_t\) and spectral-flow
\((\Delta/E)\dot\Delta\,\partial_E\) structures that one expects on the
left-hand side of \cref{eq:target}; \cref{sec:scalar_reduction} spells out
which pieces are actually generated by the plain-trace matrix projection.

The collision self-energy splits as
\(\check\Sigma=\check\Sigma_{\mathrm{imp}}+\check\Sigma_{\mathrm{inel}}\),
where the elastic impurity self-energy in the \(s\)-wave Born approximation is
\begin{equation}
  \check\Sigma_{\mathrm{imp}}(\br,t;E)
  =
  -\frac{\ii\hbar}{2\tau}\,\avg{\checkg(E,\hat\bp',\br,t)}_{\hat\bp'},
  \qquad
  \frac{1}{\tau}=\frac{2\pi}{\hbar}\, n_i N_0 u^2,
  \label{eq:sigma_imp_swave_check}
\end{equation}
isotropic in \(\hat\bp\) and identical in form for the \(R/A/K\) blocks. Here
\(\avg{\cdot}_{\hat\bp'}\) is the angular average over the Fermi surface,
\(n_i\) the impurity concentration, \(N_0\) the single-spin normal-state
density of states at the Fermi level, and \(u\) the (constant) \(s\)-wave
scattering matrix element. The inelastic part \(\check\Sigma_{\mathrm{inel}}\)
contains electron-phonon, photon, and electron-electron contributions; their
explicit forms are not needed for the angular reduction and are carried
schematically.

\subsection{Dirty-limit small parameter and angular ansatz}
\label{sec:angular_ansatz}

\subsubsection{Hierarchy of energy scales}

The dirty limit assumes the elastic impurity rate dominates every other
energy scale relevant to the kinetic equation:
\begin{equation}
  \frac{\hbar}{\tau}
  \;\gg\;
  \Delta,\;
  \hbar\,\partial_t,\;
  \hbar\vF/L,\;
  \hbar/\tau_{\mathrm{inel}},
  \label{eq:dirty_hierarchy}
\end{equation}
where \(L\) is the spatial scale of variation of \(\Delta(\br)\) and of
\(\fL,\fT\) (i.e.\ inverse gradient length), and \(\tau_{\mathrm{inel}}\) is
any inelastic relaxation time. Equivalently \(\ell\equiv\vF\tau\ll\xi_0,L\).
This is the same condition used in Kopnin~\S 5.6, \S 9.3, and \S 10.5; it is
also the condition imposed in the qp-diffusion paper to label the resulting
operator as ``Usadel.''

\subsubsection{Angular harmonic decomposition}

Decompose the matrix propagator on the Fermi surface as
\begin{equation}
  \checkg(E,\hat\bp,\br,t)
  =
  \checkg_0(E,\br,t)
  + \hat\bp\!\cdot\!\check{\mathbf{g}}_1(E,\br,t)
  + \cdots,
  \qquad
  \checkg_0\equiv\avg{\checkg}_{\hat\bp},
  \label{eq:angular_decomp}
\end{equation}
with the higher harmonics suppressed by additional powers of the small
parameter \(\tau/L \cdot \vF\) (equivalently \(\ell/L\)). The normalization
\(\checkg\circ\checkg=\openone\) gives, to first harmonic,
\begin{equation}
  \checkg_0\circ\checkg_0=\openone,
  \qquad
  \{\checkg_0,\check{\mathbf{g}}_1\}_{\circ}=0,
  \label{eq:norm_harmonic}
\end{equation}
i.e.\ the first-harmonic correction anticommutes with the isotropic part.

\subsubsection{Where the impurity term dominates}

Substituting \cref{eq:sigma_imp_swave_check} and \cref{eq:angular_decomp} into
the impurity self-energy gives
\begin{equation}
  \check\Sigma_{\mathrm{imp}}(\br,t;E)
  =
  -\frac{\ii\hbar}{2\tau}\,\checkg_0(E,\br,t),
  \label{eq:Sigma_imp_isotropic}
\end{equation}
since the \(\hat\bp'\) average kills the first-harmonic piece. So
\(\check\Sigma_{\mathrm{imp}}\) is itself angle-independent and proportional
to \(\checkg_0\). Accordingly, the impurity commutator in
\cref{eq:keldysh_eilenberger} is
\begin{equation}
  -\bigl[\check\Sigma_{\mathrm{imp}},\,\checkg\bigr]
  =
  \frac{\ii\hbar}{2\tau}\,
  \bigl[\checkg_0,\,\checkg\bigr]
  =
  \frac{\ii\hbar}{2\tau}\,
  \hat\bp\!\cdot\!\bigl[\checkg_0,\,\check{\mathbf{g}}_1\bigr]
  + O(\hat\bp\hat\bp).
  \label{eq:imp_comm_first_harmonic}
\end{equation}
The isotropic part of this commutator vanishes (\(\checkg_0\) commutes with
itself), so the impurity term contributes \emph{only} to the first-harmonic
balance --- and there it scales as \(\hbar/\tau\), the largest scale in the
hierarchy \cref{eq:dirty_hierarchy}. This is the structural reason the
impurity collision integral cannot survive on the RHS of the reduced
equation: in the only equation it touches at leading order (the first-harmonic
equation), it dominates everything else and gets used to solve for
\(\check{\mathbf{g}}_1\).

\subsection{Outline of the derivation}
\label{sec:roadmap}

The remainder of the derivation proceeds in three steps, each of which is a
self-contained calculation that we will spell out turn by turn:
\begin{description}
  \item[\cref{sec:matrix_usadel} Matrix Keldysh--Usadel.] Take the
  zeroth and first angular moments of \cref{eq:keldysh_eilenberger}; in the
  dirty limit, the first-harmonic equation slaves
  \(\check{\mathbf{g}}_1\approx-\ell\,\checkg_0\,\widetilde\bnabla\checkg_0\),
  which substituted back into the zeroth-moment equation yields the matrix
  Keldysh--Usadel equation
  \(\hbar\DN\,\widetilde\bnabla\!\cdot\!(\checkg_0\,\widetilde\bnabla\checkg_0)
  +[\,E\tau_3-\hatD-\check\Sigma_{\mathrm{inel}},\,\checkg_0\,]_{\circ}=0\)
  with \(\checkg_0\circ\checkg_0=\openone\).
  \item[\cref{sec:fLfT} Coupled \((\fL,\fT)\) kinetic equations.] Project the
  Keldysh component using
  \cref{eq:dist_matrix_ansatz}. The \(\Tr\) and \(\Tr[\tau_3\,\cdot]\)
  projections give two decoupled scalar equations with diffusion-type spatial
  operators carrying spectral coefficients
  \(\mathcal{D}_{L/T}=\tfrac{1}{4}\Tr[\openone-\hatg^{R}(\tau_3)\hatg^{A}(\tau_3)]\),
  evaluating to \(\mathcal{D}_L=1\) and
  \(\mathcal{D}_T=\DOS^2\) on the bulk BCS background above the gap.
  \item[\cref{sec:scalar_reduction} Scalar reduction.] Set
  \(\fT\!\to\!0\) (justified when charge-imbalance relaxation is fast and
  there is no charge-injection source) and translate \(\fL=1-2f\). The
  plain-trace projection produces \cref{eq:target} in conservative form,
  including the time-dependent spectral-flow term; the moving-energy-shell
  and adiabatic branch-projector constructions in
  \cref{sec:coordinate_lift,sec:moving_bin_projection,sec:branch_projector_lift}
  supply its kinematic interpretation.
\end{description}

\noindent The \emph{branch-projection} alternative --- project onto the
positive-energy Bogoliubov quasiparticle branch first, \emph{then} angular-average
with the impurity collision integral on the right-hand side --- is a
logically separate derivation and is treated in
\cref{sec:branch_alternative}. Its moment reduction carries the spectral
weight \(\DOS\) outside the divergence, so with the BRT closure for
\(\tautr(E)\) it lands on the \emph{same} \((p,q)=(1,0)\) operator as the
dirty-limit Keldysh--Usadel projection; the \((0,-1)\) and \((0,-2)\)
entries of the main-text taxonomy are legacy placements (Fick ans\"atze
for the occupation itself), not the end point of the branch reduction.
The projection/averaging non-commutation posed in
main-text~\cref{M-eq:projection_average_commutator} is exhausted by the
branch-to-trajectory relabelling of odd angular harmonics and marks a
boundary between representations, not between diffusion laws.

\subsection{Step 1: matrix Keldysh--Usadel}
\label{sec:matrix_usadel}

This step reduces \cref{eq:keldysh_eilenberger} to a matrix
diffusion equation for the angularly averaged quasiclassical propagator
\(\checkg_0\), with the elastic impurity rate absorbed into the diffusion
coefficient \(\DN=\vF^2\tau/d\). It is the Keldysh generalization of the
Matsubara reduction in Kopnin~\S 5.6. The structure is the same; only the labels (real-time \(E\) and
\(R/A/K\) blocks instead of Matsubara \(i\omega_n\)) change.

\subsubsection{Angular moments}
\label{sec:angular_moments}

Insert the angular harmonic decomposition
\cref{eq:angular_decomp} into the Keldysh--Eilenberger equation
\cref{eq:keldysh_eilenberger}. The streaming term separates by parity in
\(\hat\bp\):
\begin{equation}
  \ii\hbar\vF\,\hat\bp\!\cdot\!\widetilde\bnabla\,\checkg
  =
  \ii\hbar\vF\,\hat\bp\!\cdot\!\widetilde\bnabla\checkg_0
  +
  \ii\hbar\vF\,\hat p_i\hat p_j\,\widetilde\nabla_i\,\check g_{1,j}
  + O(\hat p\hat p\hat p),
  \label{eq:streaming_split}
\end{equation}
and the spatially isotropic operators in the commutator commute with the
angular average:
\begin{equation}
  [\,E\tau_3-\hatD-\check\Sigma_{\mathrm{inel}},\,\checkg\,]_{\circ}
  =
  [\,E\tau_3-\hatD-\check\Sigma_{\mathrm{inel}},\,\checkg_0\,]_{\circ}
  + \hat\bp\!\cdot\![\,E\tau_3-\hatD-\check\Sigma_{\mathrm{inel}},\,\check{\mathbf g}_1\,]_{\circ},
  \label{eq:iso_comm_split}
\end{equation}
where we have assumed the inelastic self-energy is isotropic in \(\hat\bp\)
to leading order (true for \(s\)-wave electron-phonon scattering with
isotropic Eliashberg coupling and for the standard
relaxation-time-approximation forms; the validity of dropping the
first-harmonic part of \(\check\Sigma_{\mathrm{inel}}\) is the same condition
that lets phonon scattering be treated as a single-rate process in the
homogeneous Fischer--Catelani kinetic equations). The elastic impurity
contribution from \cref{eq:imp_comm_first_harmonic} sits entirely in the
first-harmonic sector.

Take the zeroth angular moment,
\(\langle\cdot\rangle_{\hat\bp}\), of \cref{eq:keldysh_eilenberger}:
\begin{equation}
  \boxed{\,
  \frac{\ii\hbar\vF}{d}\,\widetilde\bnabla\!\cdot\!\check{\mathbf g}_1
  +
  \bigl[\,E\tau_3-\hatD-\check\Sigma_{\mathrm{inel}},\,\checkg_0\,\bigr]_{\circ}
  =
  0,
  \,}
  \label{eq:moment0}
\end{equation}
using \(\langle\hat p_i\hat p_j\rangle_{\hat\bp}=\delta_{ij}/d\) on the
streaming and \(\langle\hat\bp\rangle=0\) on every \(\hat\bp\)-linear piece,
including the impurity commutator
\cref{eq:imp_comm_first_harmonic}. Take the first angular moment,
\(d\,\langle\hat\bp\,\cdot\rangle_{\hat\bp}\) (the factor \(d\) recovers
\(\check{\mathbf g}_1=d\,\langle\hat\bp\checkg\rangle_{\hat\bp}\)):
\begin{equation}
  \boxed{\,
  \ii\hbar\vF\,\widetilde\bnabla\,\checkg_0
  +
  \bigl[\,E\tau_3-\hatD-\check\Sigma_{\mathrm{inel}},\,\check{\mathbf g}_1\,\bigr]_{\circ}
  +
  \frac{\ii\hbar}{2\tau}\,\bigl[\checkg_0,\,\check{\mathbf g}_1\bigr]
  =
  0,
  \,}
  \label{eq:moment1}
\end{equation}
where the impurity contribution comes from
\cref{eq:imp_comm_first_harmonic}: with
\(\check\Sigma_{\mathrm{imp}}=-(\ii\hbar/2\tau)\langle\checkg\rangle\) and the
identity
\(-[\check\Sigma_{\mathrm{imp}},\checkg_0+\hat\bp\cdot\check{\mathbf g}_1]
=(\ii\hbar/2\tau)\,\hat\bp\cdot[\checkg_0,\check{\mathbf g}_1]\),
the \(d\langle\hat\bp\,\cdot\rangle\) projection produces
\((\ii\hbar/2\tau)[\checkg_0,\check{\mathbf g}_1]\) on the left-hand side of
\cref{eq:moment1}.

\subsubsection{Dirty-limit slaving of the first harmonic}
\label{sec:slaving}

In \cref{eq:moment1} the impurity commutator (size \(\hbar/\tau\)) dominates
the spectral and inelastic commutators with \(\check{\mathbf g}_1\) (sizes
\(E,\Delta,\hbar/\tau_{\mathrm{inel}}\)) and the Moyal corrections of order
\(\hbar\partial_t\) by the dirty hierarchy \cref{eq:dirty_hierarchy}. It does
\emph{not} dominate the streaming term itself: with
\(\check{\mathbf g}_1\sim(\ell/L)\,\checkg_0\), both the streaming
\(\ii\hbar\vF\,\widetilde\bnabla\checkg_0\sim\hbar\vF/L\) and the impurity
commutator
\((\ii\hbar/2\tau)[\checkg_0,\check{\mathbf g}_1]\sim(\hbar/\tau)(\ell/L)
=\hbar\vF/L\) are of the same order. The dirty limit \emph{balances} these
two against each other, with the spectral and inelastic commutators with
\(\check{\mathbf g}_1\) appearing only as subleading corrections:
\begin{equation}
  \ii\hbar\vF\,\widetilde\bnabla\,\checkg_0
  +
  \frac{\ii\hbar}{2\tau}\,\bigl[\checkg_0,\,\check{\mathbf g}_1\bigr]
  \approx
  0,
  \qquad
  \Longleftrightarrow
  \qquad
  \vF\,\widetilde\bnabla\,\checkg_0
  \approx
  -\frac{1}{2\tau}\,\bigl[\checkg_0,\,\check{\mathbf g}_1\bigr].
  \label{eq:dirty_balance}
\end{equation}
The factor \(\ii\hbar\) cancels between the two terms; the slaving relation
is therefore real, as it must be for the slaved \(\check{\mathbf g}_1\) to be
consistent with a real diffusion current.
Solve for \(\check{\mathbf g}_1\). Left-multiply by \(\checkg_0\) and use
the leading-gradient simplifications: at the order at which we have dropped
the spectral and inelastic commutators with \(\check{\mathbf g}_1\), the
Moyal product reduces to ordinary matrix multiplication, so
\(\checkg_0\circ\checkg_0=1\) becomes \(\checkg_0^2=1\) and the
first-harmonic anticommutation \(\{\checkg_0,\check{\mathbf g}_1\}=0\) from
\cref{eq:norm_harmonic} becomes the ordinary matrix anticommutator. Then
\begin{equation}
  \checkg_0\,\bigl[\checkg_0,\check{\mathbf g}_1\bigr]
  =
  \checkg_0\checkg_0\,\check{\mathbf g}_1
  -
  \checkg_0\,\check{\mathbf g}_1\,\checkg_0
  =
  \check{\mathbf g}_1
  -
  (-\check{\mathbf g}_1\checkg_0)\checkg_0
  =
  \check{\mathbf g}_1
  +
  \check{\mathbf g}_1\,\checkg_0^2
  =
  2\,\check{\mathbf g}_1,
  \label{eq:trick_identity}
\end{equation}
the identity that makes the Usadel reduction work cleanly. Hence
\begin{equation}
  \checkg_0\,\vF\,\widetilde\bnabla\,\checkg_0
  \approx
  -\frac{1}{\tau}\,\check{\mathbf g}_1,
\end{equation}
giving the first-harmonic slaving
\begin{equation}
  \boxed{\,
  \check{\mathbf g}_1
  \;\approx\;
  -\,\ell\,\checkg_0\,\widetilde\bnabla\,\checkg_0,
  \qquad
  \ell=\vF\tau.
  \,}
  \label{eq:slaving}
\end{equation}
This is the Keldysh-matrix version of the standard Usadel slaving (Kopnin
Eq.~5.96). The slaving is uniform in the
Keldysh structure: the same expression holds with \(\checkg_0\) replaced by
\(\hatg_0^R\), \(\hatg_0^A\), or by the full triangular matrix \(\checkg_0\)
including the \(K\) block, because the only inputs are the impurity
self-energy form, the normalization, and the anticommutator
\(\{\checkg_0,\check{\mathbf g}_1\}=0\) --- all of which respect the Keldysh
block structure.

\paragraph{Ingredients.} The factor of two in
\cref{eq:trick_identity} is the only nontrivial algebraic step. It uses
exactly two facts: (i) \(\checkg_0^2=1\), which is the leading-order content
of the quasiclassical normalization, and (ii)
\(\{\checkg_0,\check{\mathbf g}_1\}=0\), which is the first-harmonic
content of the same normalization. Both facts already presuppose the
gradient/dirty expansion (we kept only the leading \(O(\hat\bp^0)\) and
\(O(\hat\bp^1)\) terms in \cref{eq:norm_harmonic}). Higher-harmonic
corrections to \cref{eq:slaving} appear at order
\(\bigl(\ell/L\bigr)^2\), where \(L\) is the spatial scale of variation; we
treat these as negligible in the dirty regime.

\subsubsection{Matrix Keldysh--Usadel equation}
\label{sec:matrix_KU}

Substitute the slaving \cref{eq:slaving} into the zeroth-moment equation
\cref{eq:moment0}:
\begin{equation}
  \frac{\ii\hbar\vF}{d}\,\widetilde\bnabla\!\cdot\!
  \bigl(-\ell\,\checkg_0\,\widetilde\bnabla\,\checkg_0\bigr)
  +
  \bigl[E\tau_3-\hatD-\check\Sigma_{\mathrm{inel}},\,\checkg_0\bigr]_{\circ}
  =
  0,
\end{equation}
which collapses, using \(\DN=\vF\ell/d=\vF^2\tau/d\) and multiplying through
by \(\ii\), to the matrix Keldysh--Usadel equation
\begin{equation}
  \boxed{\,
  \hbar\DN\,\widetilde\bnabla\!\cdot\!
  \bigl(\checkg_0\,\widetilde\bnabla\,\checkg_0\bigr)
  +
  \ii\,\bigl[E\tau_3-\hatD-\check\Sigma_{\mathrm{inel}},\,\checkg_0\bigr]_{\circ}
  =
  0,
  \qquad
  \checkg_0\circ\checkg_0=1.
  \,}
  \label{eq:matrix_keldysh_usadel}
\end{equation}
This form is equivalent, up to choices in the gap-matrix convention
(\(\hatD=-\ii\Delta\tau_2\) here vs.\ \(\hatD=\Delta(\ii\tau_2)\) elsewhere) and
to overall sign of the equation, to the matrix Keldysh--Usadel equation
written in the Belzig--Kopnin--Wilhelm review~\cite{belzig1999} and in
Kopnin's \S 9.3--10.5. We will not insist on identity with any particular
printed form below \cref{eq:matrix_keldysh_usadel}; what matters is that
this is one self-consistent derivation. The factor of \(\ii\) on the spectral
commutator is the place where the time-derivative \(\hbar\partial_t\) term
appears on the left-hand side of the kinetic equation when we expand the Moyal
commutator in \cref{sec:fLfT,sec:scalar_reduction}; whether that expansion
also generates the spectral-flow \((\Delta/E)\dot\Delta\,\partial_E\) term is
an open point examined in \cref{sec:scalar_reduction}. We defer the explicit
calculation of those factors and signs to that derivation rather than
asserting them here.

\paragraph{What this equation contains.} \cref{eq:matrix_keldysh_usadel}
is a single matrix equation for \(\checkg_0(E,\br,t)\), with all of the
information about elastic impurity scattering compressed into the diffusion
constant \(\DN=\vF^2\tau/d\). Its retarded and advanced components alone are
the spectral Usadel equations. Its
Keldysh component, parameterized through the distribution matrix
\cref{eq:dist_matrix_ansatz}, is the matrix kinetic equation we will project
onto the longitudinal and transverse scalar modes in
\cref{sec:fLfT}. The Moyal commutator \([\cdot,\cdot]_{\circ}\) still
contains the Wigner gradient corrections in \((E,t)\) and \((\bp,\br)\); for
the Keldysh component these corrections produce both the
\(\partial_t\) term and the
spectral-flow term \((\Delta/E)\dot\Delta\partial_E\) directly
(\cref{sec:scalar_reduction}; the moving-shell construction and the adiabatic
retarded spectral equation \([\hat L_0,\hatg^R]_\star=0\) supply its
kinematic reading, and \cref{sec:resolution} confirms that a moving-frame
reformulation adds nothing).

\subsection{Step 2: Keldysh projection and the spatial flux dressing}
\label{sec:fLfT}

In this section we take the Keldysh component of the matrix
Keldysh--Usadel equation \cref{eq:matrix_keldysh_usadel}, substitute the
distribution-matrix ansatz
\(\hat h=\fL\openone+\fT\tau_3\), and compute the spatial flux that
appears on the left-hand side of the kinetic equation. The result is the
undressed coefficient \(\mathcal D_L=1\) for the longitudinal mode (above
the gap; \(0\) below) and the dressing
\(\mathcal D_T=\DOS^2\) for the transverse mode --- i.e.\ the exponent
\(q=0\) on the longitudinal flux of the main text, with the squared-DOS
dressing residing in the transverse mode. The complementary calculation for
the time-derivative side, which fixes the time-weight exponent \(p\) and
produces the time-dependent spectral-flow term, is carried out in
\cref{sec:scalar_reduction}.

\subsubsection{Keldysh component of the matrix Usadel equation}
\label{sec:K_block}

The triangular Keldysh block structure of \cref{eq:matrix_keldysh_usadel}
implies, for any product
\(\check A\circ\check B\),
\((\check A\circ\check B)^K=\hat A^R\circ\hat B^K+\hat A^K\circ\hat B^A\)
and analogous diagonal rules for the \(R\) and \(A\) blocks. Separate the
coherent spectral operator from the collision self-energies: write
\(\hat L_0\equiv E\tau_3-\hatD\), which has no \(K\) component (it is a
``classical'' operator in Keldysh space, equal on R and A blocks), and
collect all retarded/advanced/Keldysh contributions of the inelastic
self-energy \(\check\Sigma_{\mathrm{inel}}\) into a single
collision-integral object
\begin{equation}
  \Icoll[\checkg_0]
  \equiv
  \hatg_0^{R}\circ\hat\Sigma_{\mathrm{inel}}^{K}
  +\hatg_0^{K}\circ\hat\Sigma_{\mathrm{inel}}^{A}
  -\hat\Sigma_{\mathrm{inel}}^{R}\circ\hatg_0^{K}
  -\hat\Sigma_{\mathrm{inel}}^{K}\circ\hatg_0^{A},
  \label{eq:Icoll_def}
\end{equation}
which is the full \(K\)-block of
\([-\check\Sigma_{\mathrm{inel}},\checkg_0]_{\circ}\) (with the minus sign
absorbed into the definition so that
\(\Icoll\) appears with a \(+\) sign in the K-block kinetic equation
below). It reduces to the standard electron--phonon, photon, and
electron--electron collision integrals on substituting the appropriate
self-energies (Eliashberg~\cite{eliashberg1972}; Kopnin~\S 10.4). With this bookkeeping the
\(K\)-block of \cref{eq:matrix_keldysh_usadel} reads
\begin{equation}
  \boxed{\,
  \hbar\DN\,\widetilde\bnabla\!\cdot\!
  \bigl(\checkg_0\,\widetilde\bnabla\,\checkg_0\bigr)^{K}
  +
  \ii\,\bigl[\,\hat L_0,\,\hatg_0^{K}\,\bigr]_{\circ}
  +
  \ii\,\Icoll[\checkg_0]
  =
  0,
  \,}
  \label{eq:K_block_kinetic}
\end{equation}
a single matrix equation in Nambu space for the Keldysh-block propagator
\(\hatg_0^{K}(E,\br,t)\). The two physical scalar modes \(\fL,\fT\) are
extracted from \(\hatg_0^{K}\) by the distribution-matrix parameterization
introduced next; the projection that selects each mode is what determines
the spectral coefficients on both sides.

\subsubsection{Distribution-matrix ansatz and the BCS spectral functions}
\label{sec:bcs_spectral}

Insert the parameterization \cref{eq:dist_matrix_ansatz} into the K-block
\cref{eq:K_block_kinetic}. For BCS bulk above the gap, the retarded and
advanced quasiclassical propagators are (main text label
main-text~\cref{M-eq:bcs_advanced_matrix})
\begin{equation}
  \hatg^{R}
  =
  \DOS\,\tau_3 + \ii N_2\,\tau_2,
  \qquad
  \hatg^{A}
  =
  -\DOS\,\tau_3 - \ii N_2\,\tau_2,
  \label{eq:gRA_BCS}
\end{equation}
in the parent-paper basis \(\{\tau_3,\ii\tau_2\}\), where
\(\DOS,N_2\) are real and positive on the quasiparticle branch
\(E>\Delta\). The advanced propagator is the conjugation image
\(\hatg^A=-\tau_3\hatg^{R\dagger}\tau_3\), which above the gap collapses to
\(\hatg^A=-\hatg^R\) (main text, \cref{M-eq:bcs_advanced_matrix}); the
normalization \((\hatg^{R/A})^2=\openone\) follows from the BCS spectral
identity
\[
  \DOS^2-N_2^2=1
\]
but does not by itself fix the relative sign---that is fixed by the
spectral equation and the equilibrium Keldysh anomalous weight.
Direct substitution and the Pauli algebra
\(\tau_a\tau_b=\delta_{ab}\openone+\ii\epsilon_{abc}\tau_c\) give the bilinears
\begin{align}
  \hatg^{R}\hatg^{A}
  &=
  -\bigl(\DOS^{2}-N_2^{2}\bigr)\,\openone
  =
  -\openone,
  \label{eq:gRgA}
  \\
  \hatg^{R}\tau_3\hatg^{A}
  &=
  -(\DOS^{2}+N_2^{2})\,\tau_3
  \;-\; 2\ii\,\DOS\,N_2\,\tau_2 .
  \label{eq:gRtau3gA}
\end{align}
The first line is exact at the matrix level: with \(\hatg^A=-\hatg^R\)
above the gap, \(\hatg^R\hatg^A=-(\hatg^R)^2=-\openone\), so the
longitudinal coefficient matrix \(\openone-\hatg^R\hatg^A=2\,\openone\) is
proportional to the identity and carries \emph{no} off-diagonal remnant.
The cross-term structure instead appears in the second line (equivalently
in \(\hatg^R\tau_3\hatg^A\tau_3=-(\DOS^2+N_2^2)\openone+2\DOS N_2\tau_1\)):
the \(\tau_1\) piece survives at the matrix level in the transverse
coefficient, drops out after the Nambu trace projection in
\cref{sec:DL_DT}, but is still present in the matrix coefficient of
\(\widetilde\bnabla\fT\) inside the divergence of the matrix flux
\cref{eq:matrix_current_general}, and any reader doing the calculation
directly should keep it. These two evaluations are the load-bearing inputs
to the longitudinal and transverse flux coefficients computed next.

The Keldysh component of the angularly averaged propagator becomes,
substituting \(h=\fL\openone+\fT\tau_3\) and
\((\hatg^{R}-\hatg^{A})=2\DOS\tau_3+2\ii N_2\tau_2\),
\(\hatg^{R}\tau_3-\tau_3\hatg^{A}=2\DOS\,\openone\),
\begin{equation}
  \hatg_0^{K}
  =
  \hatg^{R}\hat h - \hat h\,\hatg^{A}
  =
  (2\DOS\,\tau_3+2\ii N_2\,\tau_2)\,\fL
  + 2\DOS\,\fT\,\openone .
  \label{eq:gK_BCS_h}
\end{equation}
The trace identities of the main text follow immediately:
\begin{equation}
  \Tr[\hatg_0^{K}] = 4\,\DOS\,\fT,
  \qquad
  \Tr[\tau_3\,\hatg_0^{K}] = 4\,\DOS\,\fL,
  \label{eq:trace_identities}
\end{equation}
showing that \emph{plain trace} of \(\hatg_0^K\) projects onto the
transverse mode \(\fT\) (with weight \(\DOS\)), while the
\emph{\(\tau_3\)-trace} of \(\hatg_0^K\) projects onto the longitudinal
mode \(\fL\) (with weight \(\DOS\)).

\subsubsection{Spatial flux: matrix-current decomposition}
\label{sec:matrix_current}

Compute the K-block of the matrix flux
\((\checkg_0\widetilde\bnabla\checkg_0)^{K}\) using the parameterization
\(\hatg_0^K=\hatg^{R}\hat h-\hat h\hatg^{A}\). Expanding
\(\widetilde\bnabla\hatg_0^K\) and using \((\hatg^{R/A})^2=\openone\)
gives, with \(\hat h=\fL\openone+\fT\tau_3\) (so the \(\widetilde\bnabla\hat h\)
factors \(\widetilde\bnabla\fL\openone+\widetilde\bnabla\fT\tau_3\) are c-number
coefficients on Pauli matrices):
\begin{equation}
  \begin{aligned}
  (\checkg_0\,\widetilde\bnabla\,\checkg_0)^{K}
  =\;&
  (\openone-\hatg^{R}\hatg^{A})\,\widetilde\bnabla\fL
  +(\tau_3-\hatg^{R}\tau_3\hatg^{A})\,\widetilde\bnabla\fT
  \\
  &
  +\hatg^{R}\widetilde\bnabla\hatg^{R}\cdot\hat h
  -\hat h\cdot\hatg^{A}\widetilde\bnabla\hatg^{A}.
  \end{aligned}
  \label{eq:matrix_current_general}
\end{equation}
The first line is the diffusive-flux content, with the spatial gradients
of \(\fL,\fT\) appearing explicitly. The second line is the
spectral-anomalous drift, with \(\hat h\) on the left of the
\(\hatg^{A}\widetilde\bnabla\hatg^{A}\) factor and on the right of the
\(\hatg^{R}\widetilde\bnabla\hatg^{R}\) factor. For
spatially uniform \(\Delta(\br)\) the drift vanishes identically because
\(\widetilde\bnabla\hatg^{R/A}=0\). For varying real \(\Delta(\br)\) it
collapses as well: with
\(\hatg^A=-\hatg^R\) above the gap,
\(\hatg^A\widetilde\bnabla\hatg^A=\hatg^R\widetilde\bnabla\hatg^R\), so the
drift reduces to the commutator
\(\fT\,[\hatg^R\widetilde\bnabla\hatg^R,\tau_3]\). Since
\(\hatg^R\widetilde\bnabla\hatg^R=(\DOS\widetilde\bnabla N_2-N_2\widetilde\bnabla\DOS)\tau_1\)
for the real-gap BCS matrices, the commutator is \(\propto\tau_2\) and both
its Nambu trace and its \(\tau_3\)-trace vanish: there is \emph{no}
DOS-gradient drift and no \(\fL\)--\(\fT\) mixing in either projected
channel (symbolically verified identities I7/I8). The general (including non-BCS)
spatially varying spectrum is carried out in
\cref{app:proximity}.

For BCS bulk, substitute \cref{eq:gRgA} and \cref{eq:gRtau3gA}:
\begin{equation}
  \openone-\hatg^{R}\hatg^{A}
  =
  2\,\openone,
  \qquad
  \tau_3-\hatg^{R}\tau_3\hatg^{A}
  =
  2\,\DOS^{2}\,\tau_3
  \;+\; 2\ii\,\DOS\,N_2\,\tau_2,
  \label{eq:flux_dressings_BCS}
\end{equation}
where the transverse simplification used \(1+\DOS^{2}+N_2^{2}=2\DOS^{2}\)
on the \(\tau_3\) coefficient (from \(\DOS^2-N_2^2=1\)) and retained the
\(\tau_2\) cross-term from \cref{eq:gRtau3gA}. Substituting back into
\cref{eq:matrix_current_general} for uniform \(\Delta\) (where the
spectral-anomalous drift in the second line of \cref{eq:matrix_current_general}
vanishes):
\begin{equation}
  (\checkg_0\,\widetilde\bnabla\,\checkg_0)^{K}
  =
  2\,\openone\,\widetilde\bnabla\fL
  + \bigl(2\,\DOS^{2}\,\tau_3+2\ii\,\DOS\,N_2\,\tau_2\bigr)\,\widetilde\bnabla\fT.
  \label{eq:matrix_current_BCS}
\end{equation}
The matrix coefficient of \(\widetilde\bnabla\fL\) is clean: a single
\(2\,\openone\) piece, already proportional to the identity at the matrix
level. The matrix coefficient of \(\widetilde\bnabla\fT\) carries
\emph{both} a diagonal \(2\DOS^{2}\tau_3\) piece and an off-diagonal
\(2\ii\DOS N_2\tau_2\) piece; the \(\tau_2\) piece survives at the matrix
level and only drops out under the \(\tau_3\)-trace projection in
\cref{sec:DL_DT}.

\subsubsection{Spectral coefficients D\_L and D\_T}
\label{sec:DL_DT}

Following Belzig et~al.~\cite{belzig1999} (their Eqs.~38 and~37,
respectively), define the
longitudinal and transverse spectral coefficients
\begin{equation}
  \mathcal D_L(E,\br)
  \equiv
  \tfrac{1}{4}\Tr\!\bigl[\openone-\hatg^{R}\hatg^{A}\bigr],
  \qquad
  \mathcal D_T(E,\br)
  \equiv
  \tfrac{1}{4}\Tr\!\bigl[\openone-\hatg^{R}\tau_3\hatg^{A}\tau_3\bigr],
  \label{eq:DL_DT_def}
\end{equation}
which are the BCS-bulk evaluations of the Nambu traces of the matrix
coefficients of \(\widetilde\bnabla\fL\) and (after \(\tau_3\)-conjugation)
\(\widetilde\bnabla\fT\). Plugging in \cref{eq:gRgA}, \cref{eq:gRtau3gA}
and using \(\Tr[\openone]=2\), \(\Tr[\tau_a]=0\):
\begin{equation}
  \mathcal D_L
  =
  \tfrac{1}{4}\Tr\bigl[2\,\openone\bigr]
  =
  1,
  \qquad
  \mathcal D_T
  =
  \tfrac{1}{4}\Tr\bigl[2\DOS^{2}\openone-2\DOS N_2\tau_1\bigr]
  =
  \DOS^{2}(E,\br),
  \label{eq:DL_DT_BCS}
\end{equation}
matching main-text~\cref{M-eq:DL_DT_BCS_bulk}. The \(\tau_1\)
piece in the transverse coefficient
(\(\openone-\hatg^R\tau_3\hatg^A\tau_3=2\DOS^2\openone-2\DOS N_2\tau_1\))
is killed by the trace; this is where the matrix-level cross-term becomes
inconsequential, even though it was present in the matrix coefficient
itself.

These coefficients are the projected scalar spectral currents for the
longitudinal and transverse modes,
\begin{equation}
  \mathbf j_L
  =
  -\DN\,\mathcal D_L(E,\br)\,\widetilde\bnabla\fL
  =
  -\DN\,\widetilde\bnabla\fL,
  \label{eq:jL_BCS}
\end{equation}
\begin{equation}
  \mathbf j_T
  =
  -\DN\,\mathcal D_T(E,\br)\,\widetilde\bnabla\fT
  =
  -\DN\,\DOS^{2}(E,\br)\,\widetilde\bnabla\fT,
  \label{eq:jT_BCS}
\end{equation}
i.e., the c-number scalar coefficients that survive after taking the
appropriate Nambu trace projection of the matrix flux. The raw matrix
flux \cref{eq:matrix_current_BCS} retains the \(\tau_2\) cross-term and
the \(\tau_3\) and \(\openone\) pieces side by side; \cref{eq:jL_BCS},
\cref{eq:jT_BCS} are what those reduce to after projecting onto the
\(\fL\) and \(\fT\) sectors with \(\Tr\) and \(\Tr[\tau_3\,\cdot]\)
respectively. The longitudinal current is \emph{undressed}, with the bare
normal-state diffusion coefficient \(\DN\) above the gap (and blocked,
\(\mathcal D_L=0\), below it); the transverse current carries the enhanced
(singular at the gap edge) spectral coefficient \(\DOS^{2}(E)\). This
asymmetry is the matrix-Usadel statement of the ``flux-dressing'' exponent
\(q=0\) for the longitudinal mode in the main text's \(L_{p,q}\)
classification (main-text~\cref{M-eq:Lpq_family}).
The exponent \(q=0\) is fixed; the time-weight exponent \(p\) is fixed in
\cref{sec:scalar_reduction}.

\paragraph{Ingredients.} The spatial-flux assignment follows entirely
from (i) the BCS bilinears
\(\hatg^{R}\hatg^{A}=-\openone\) and
\(\hatg^{R}\tau_3\hatg^{A}=-(\DOS^2+N_2^2)\tau_3-2\ii\DOS N_2\tau_2\),
which evaluate the matrix
coefficients of \(\widetilde\bnabla\fL\) and \(\widetilde\bnabla\fT\) on
the BCS bulk background; (ii) the BCS sum rule \(\DOS^2-N_2^2=1\), which
collapses \(\hatg^{R}\hatg^{A}\) to \(-\openone\) at the matrix level and
replaces \(1+N_2^2\) with \(\DOS^2\) in the transverse trace;
and (iii) the matrix-current decomposition
\cref{eq:matrix_current_general}, which separates the diffusive-flux
pieces from the spectral-anomalous drift. The longitudinal coefficient is
proportional to the identity \emph{before} any trace is taken, so
the conclusion \(\mathcal D_L=1\) is robust. None of these uses the
time-derivative side, the inelastic self-energy structure, or any choice
of conserved density. The undressed longitudinal flux
\(\mathcal D_L=1\) (above the gap) is therefore unambiguous in the dirty
matrix-Usadel route, regardless of how the time weight is fixed in
\cref{sec:scalar_reduction}.

\paragraph{Bulk vs.\ proximity caveat.} The simplifications
\cref{eq:gRgA}, \cref{eq:gRtau3gA} use the bulk BCS spectral functions
\(\DOS,N_2\) for the parent-paper basis
\(\hatg^{R}=\DOS\tau_3+\ii N_2\tau_2\), \(\hatg^A=-\hatg^R\). For a
general complex spectral angle \(\theta(\br,E)\) (proximity minigap,
depairing, sub-gap states) the coefficients become
\(\mathcal D_L=\cos^2(\Im\theta)\) and
\(\mathcal D_T=\cosh^2(\Re\theta)\), with the exact identity
\(\mathcal D_L\mathcal D_T=\DOS^2\); the BCS bulk values \(1\) and
\(\DOS^2\) are the real-\(\theta\) limits. This matches the dirty-limit
coefficient definitions of Belzig et~al.~\cite{belzig1999} (their
Eqs.~36--39), with the sub-gap blocking \(\mathcal D_L\to0\) following
from their Eq.~(38) evaluated on the sub-gap spectrum. The general case
is worked out in \cref{app:proximity}.

\subsection{Step 3: time-weight projection and scalar reduction}
\label{sec:scalar_reduction}

This section completes the dirty-limit reduction by computing the
time-derivative side of the K-block kinetic equation
\cref{eq:K_block_kinetic} and combining it with the spatial flux from
\cref{sec:fLfT}. The Moyal expansion of the spectral commutator
\(\ii[\hat L_0,\hat g_0^{K}]_{\circ}\) generates the time derivative
\(\hbar\partial_t\) and, for a time-dependent gap, the spectral-flow
energy flux. Plain trace projection produces a kinetic equation
for the conserved scalar density \(\DOS\fL\) (i.e., the time-weight
exponent \(p=1\)) in the canonical conservative form
\(\partial_t(\DOS\fL)+\partial_E[(\Delta/E)\dot\Delta\,\DOS\fL]\); the
scalar reduction \(\fT\to 0\) and the change of
variable \(f=(1-\fL)/2\) then yield the A1 equation. The moving-shell and
branch-projector constructions below supply the kinematic
interpretation of that result.

\subsubsection{Matrix Moyal commutator with the spectral operator}
\label{sec:moyal_spectral}

\paragraph{Spectral-operator gauge check.} The gap gauge of
\cref{sec:conventions} is fixed by one algebraic requirement: the coherent
operator \(\hat L_0=E\tau_3-\hatD=E\tau_3+\ii\Delta\tau_2\) must have the
static BCS propagator as its spectral solution. With \cref{eq:gRA_BCS},
\(\hatg^{R}=\DOS\tau_3+\ii N_2\tau_2\), and \(\DOS=E/W\), \(N_2=\Delta/W\),
\(W=\sqrt{E^2-\Delta^2}\),
\begin{equation}
  \hatg^{R}=\frac{1}{W}\bigl(E\tau_3+\ii\Delta\tau_2\bigr)=\frac{1}{W}\,\hat L_0,
  \label{eq:gR_propto_L0}
\end{equation}
so \(\hatg^{R}\propto\hat L_0\), \([\hat L_0,\hatg^{R}]=0\), and
\((\hatg^{R})^2=(\DOS^2-N_2^2)\openone=\openone\) (using
\(\{\tau_3,\tau_2\}=0\) and \(\DOS^2-N_2^2=1\)). The normalization
\((\hatg^{R})^2=\openone\) is a property of the propagator alone and holds in
any gap gauge; the two gauge-sensitive relations \(\hatg^{R}\propto\hat L_0\)
and \([\hat L_0,\hatg^{R}]=0\) fail for the alternative \(\hatD=\Delta\tau_1\),
which is why that choice is not used. The advanced propagator
\(\hatg^A=-\hatg^R\) of \cref{eq:gRA_BCS} is likewise proportional to
\(\hat L_0\), so \([\hat L_0,\hatg^{A}]=0\): retarded and advanced
propagators solve the \emph{same} static spectral equation, as they must.
(The opposite sign choice \(-\hatg^{R\dagger}\) fails this requirement,
giving \([\hat L_0,\hatg^A]=4EN_2\tau_1\neq0\); together with the
gap-equation check of main-text~\cref{M-eq:bcs_advanced_matrix}, this is what fixes
the convention.)

Adopt the Moyal product convention
\begin{equation}
  A\star B
  =
  AB + \frac{\ii\hbar}{2}\bigl(\partial_E A\cdot\partial_t B
                              -\partial_t A\cdot\partial_E B\bigr)
  + O(\hbar^2),
  \label{eq:moyal_def}
\end{equation}
where the products \(\partial_E A\cdot\partial_t B\) and
\(\partial_t A\cdot\partial_E B\) are ordinary matrix products with the
factors in the indicated left--right order. The corresponding commutator
is
\begin{equation}
  [A,B]_{\star}
  =
  [A,B]
  +\frac{\ii\hbar}{2}\bigl(\partial_E A\,\partial_t B
                          -\partial_t A\,\partial_E B
                          +\partial_t B\,\partial_E A
                          -\partial_E B\,\partial_t A\bigr)
  + O(\hbar^2).
  \label{eq:moyal_comm}
\end{equation}
The four matrix products in the parentheses do not in general combine
into a simple antisymmetric ``Poisson bracket'' for matrix-valued
operators; the ordering matters whenever \(A\) and \(B\) do not commute.

For \(A=E\tau_3\) (depending on \(E\) but not on \(t\), \(\br\), \(\bp\):
\(\partial_E A=\tau_3\), \(\partial_t A=0\)) and \(B=\hatg_0^{K}\),
\cref{eq:moyal_comm} collapses to
\begin{equation}
  [E\tau_3,\hatg_0^{K}]_{\star}
  =
  E[\tau_3,\hatg_0^{K}]
  +
  \frac{\ii\hbar}{2}\,\bigl\{\tau_3,\,\partial_t\hatg_0^{K}\bigr\},
  \label{eq:moyal_E}
\end{equation}
where \(\{\cdot,\cdot\}\) is the matrix anticommutator. Multiplying by
\(\ii\) (as the K-block kinetic equation \cref{eq:K_block_kinetic}
requires):
\begin{equation}
  \ii[E\tau_3,\hatg_0^{K}]_{\star}
  =
  \ii E[\tau_3,\hatg_0^{K}]
  -
  \frac{\hbar}{2}\,\bigl\{\tau_3,\,\partial_t\hatg_0^{K}\bigr\}.
  \label{eq:i_moyal_E}
\end{equation}
The first term is purely off-diagonal in Nambu space (a commutator with
\(\tau_3\) only), the second is the time-derivative correction. For the
time-dependent gap commutator \(A=-\hatD=\ii\Delta(t)\tau_2\)
(\(\partial_E A=0\), \(\partial_t A=\ii\dot\Delta\tau_2\)),
\cref{eq:moyal_comm} similarly gives
\begin{equation}
  \ii[\ii\Delta\tau_2,\hatg_0^{K}]_{\star}
  =
  -\Delta[\tau_2,\hatg_0^{K}]
  +
  \frac{\ii\hbar}{2}\,\dot\Delta\,\bigl\{\tau_2,\,\partial_E\hatg_0^{K}\bigr\}.
  \label{eq:i_moyal_Delta}
\end{equation}
The \(\dot\Delta\partial_E\) structure is the matrix-Usadel form of the
spectral flow: a time-dependent gap couples through the energy
derivative of \(\hatg_0^K\), with prefactor \(\dot\Delta\tau_2\).

\subsubsection{Plain-trace projection: conserved spectral density
\texorpdfstring{\(\DOS\fL\)}{N1 fL}}
\label{sec:trace_projection}

Apply the plain Nambu trace \(\Tr[\,\cdot\,]\) to the K-block kinetic
equation \cref{eq:K_block_kinetic}. Each piece projects as follows.

\paragraph{Spatial flux.} From \cref{eq:matrix_current_BCS}, the plain
trace of the BCS-bulk matrix flux at uniform \(\Delta\) is
\(\Tr[(\checkg_0\,\widetilde\bnabla\,\checkg_0)^{K}]
=4\,\widetilde\bnabla\fL\), so
\begin{equation}
  \hbar\DN\,\widetilde\bnabla\!\cdot\!
  \Tr\bigl[(\checkg_0\,\widetilde\bnabla\,\checkg_0)^{K}\bigr]
  =
  4\hbar\DN\,\widetilde\bnabla^{2}\fL :
  \label{eq:trace_diffusion}
\end{equation}
the longitudinal diffusive flux is undressed (\(\mathcal D_L=1\) above the
gap), as established in \cref{sec:DL_DT}.

\paragraph{Spectral commutator.} From \cref{eq:i_moyal_E}, the first
term \(\ii E[\tau_3,\hatg_0^{K}]\) traces to zero by cyclicity. The
second term gives
\begin{equation}
  -\frac{\hbar}{2}\Tr\bigl[\{\tau_3,\partial_t\hatg_0^{K}\}\bigr]
  =
  -\hbar\,\Tr[\tau_3\,\partial_t\hatg_0^{K}]
  =
  -4\hbar\,\partial_t(\DOS\,\fL),
  \label{eq:trace_time_deriv}
\end{equation}
using \(\Tr[\{\tau_3,X\}]=2\Tr[\tau_3 X]\) and the trace identity
\cref{eq:trace_identities}, \(\Tr[\tau_3\hatg_0^{K}]=4\DOS\fL\). The
right-hand side is the time derivative of the spectral density
\(\DOS\fL\), which is the load-bearing trace identity that fixes the
time-weight exponent \(p=1\): the matrix-Usadel projection
\emph{naturally} produces a continuity-like time derivative on the
spectral-weighted scalar \(\DOS\fL\), not on \(\fL\) alone or on
\(\DOS^2\fL\).

\paragraph{Gap commutator.} From \cref{eq:i_moyal_Delta}, the first
term \(-\Delta[\tau_2,\hatg_0^{K}]\) involves
\([\tau_2,\hatg_0^{K}]\) which, on the BCS-bulk distribution
\cref{eq:gK_BCS_h} with \(\fT=0\), evaluates to
\([\tau_2,(2\DOS\tau_3+2\ii N_2\tau_2)\fL]=2\DOS\fL[\tau_2,\tau_3]
=4\ii\DOS\fL\tau_1\),
giving \(-\Delta\cdot 4\ii\DOS\fL\tau_1=-4\ii\Delta\DOS\fL\tau_1\),
which traces to zero. The second term involves
\(\{\tau_2,\partial_E\hatg_0^{K}\}\); for \(\fT=0\),
\(\partial_E\hatg_0^{K}
=2\tau_3\,\partial_E(\DOS\fL)+2\ii\tau_2\,\partial_E(N_2\fL)\). The
\(\tau_3\) piece drops by \(\{\tau_2,\tau_3\}=0\), but the \(\tau_2\)
piece---present precisely because the corrected \(\hatg^A\) makes the
anomalous component of \(\hatg^R-\hatg^A\) nonzero---survives:
\(\{\tau_2,2\ii\tau_2\,\partial_E(N_2\fL)\}=4\ii\,\partial_E(N_2\fL)\,\openone\).
The gap commutator therefore contributes the energy-space flux
\begin{equation}
  \Tr\bigl[\ii[\ii\Delta\tau_2,\hatg_0^{K}]_{\star}\bigr]
  =
  -4\hbar\,\partial_E\!\bigl[\dot\Delta\,N_2\,\fL\bigr]
  \qquad
  (\fT=0),
  \label{eq:trace_gap}
\end{equation}
using \(\partial_E\dot\Delta=0\). Since
\(N_2=(\Delta/E)\DOS\) on the quasiparticle branch, this is exactly the
canonical spectral-flow flux
\(-4\hbar\,\partial_E[(\Delta/E)\dot\Delta\,\DOS\fL]\). For static
\(\Delta\) the contribution vanishes, so the static result is insensitive
to the gap-matrix structure; the time-dependent flux is produced in the
gap gauge of \cref{sec:conventions}, the one consistent with the
\(\tau_2\)-anomalous propagators. (For the opposite advanced-propagator
sign the \(\tau_2\) component of \(\hatg_0^K\) cancels and this flux
vanishes; the main-text conjugation relation \cref{M-eq:bcs_advanced_matrix} is what
retains it.)

Assembling the traces, the projected equation takes the canonical
conservative form
\begin{equation}
  \partial_t(\DOS\fL) + \partial_E\!\left[\tfrac{\Delta}{E}\dot\Delta\,\DOS\fL\right]
  \;=\;
  \DN\,\widetilde\bnabla^{2}\fL
  + J_1 ,
  \label{eq:expected_conservative}
\end{equation}
with the energy-space flux explicit on the left-hand side. Combined with
the BCS DOS continuity identity (labelled
\cref{M-eq:cons_dos_continuity} in the main text and reproduced here),
\begin{equation}
  \partial_t\DOS
  +
  \partial_E\!\left[\frac{\Delta}{E}\dot\Delta\,\DOS\right]
  =0,
  \label{eq:DOS_continuity}
\end{equation}
the LHS of \cref{eq:expected_conservative} collapses to the
canonical-Boltzmann form
\(\DOS\,(\partial_t+\dot E\partial_E)\fL\) with
\(\dot E\equiv(\Delta/E)\dot\Delta\), reproducing
\cref{eq:target} after \(\fL\to 1-2f\).

\paragraph{Inelastic collision integral.} The plain trace of
\(\ii\Icoll[\checkg_0]\) defines the longitudinal collision integral
\(J_1\):
\begin{equation}
  4\hbar J_1[f,n](E,\br)
  \equiv
  \ii\,\Tr\bigl[\Icoll[\checkg_0]\bigr],
  \label{eq:J1_def}
\end{equation}
where the factor \(4\hbar\) is chosen to make \(J_1\) match the
electron--phonon and photon collision integrals of Kopnin~\S 10.4; the
explicit substitution of the
inelastic self-energies is the Fermi golden-rule calculation that
produces the standard scattering, recombination, and pair-breaking
terms, and we refer to Kopnin~\S 10.4 for the calculation. \(J_1\)
depends parametrically on \(\fL,\fT\) and on the phonon distribution
\(n(\omega,t)\); the dependence on \(\fT\) drops out for an
\(s\)-wave bulk superconductor without charge injection.

The normalization deserves one explicit statement, because two
conventions coexist in the literature. Like Kopnin's sources, \(J_1\)
carries the external spectral weight \(\DOS(E)\): in the homogeneous
occupation representation, where the kinetic equation is written as
\(\partial_t f=I_{\mathrm{occ}}[f,n]\) with the Kaplan-form kernels
\cite{kaplan1976,fischer2023} (which contain only the \emph{partner}
density of states), the correspondence implied by \(\fL=1-2f\) is
\begin{equation}
  J_1(E) = -2\,\DOS(E)\,I_{\mathrm{occ}}[f,n](E).
  \label{eq:J1_occ_bridge}
\end{equation}
Occupation-form kernels must not be substituted for \(J_1\) without the
external factor \(-2\DOS(E)\).

\paragraph{What the plain trace produces.} The plain-trace projection of
the matrix Keldysh--Usadel equation produces the complete canonical
left-hand side: the \(E\tau_3\) Moyal correction gives the time
derivative \(\partial_t(\DOS\fL)\) (\cref{eq:trace_time_deriv}), and the
gap commutator's Moyal correction gives the energy-space flux
\(\partial_E[(\Delta/E)\dot\Delta\,\DOS\fL]\) (\cref{eq:trace_gap}).
Together they assemble \cref{eq:expected_conservative}, the canonical
Boltzmann form required by Kopnin~\S 10.5, directly at fixed \(E\). Three
consistency questions remain:
(i)~whether higher-order Moyal corrections to the spectral commutator
alter the leading flux,
(ii)~whether time-Moyal corrections to the spectral-anomalous drift in the
matrix current \cref{eq:matrix_current_general} contribute, and
(iii)~whether the first Wigner correction to the Keldysh ansatz shifts the
projection. The next subsection checks each explicitly; none modifies the
leading \(O(\hbar\dot\Delta)\) result.

\subsubsection{Consistency checks on the spectral-flow flux}
\label{sec:tdep_candidate_tests}

The energy-space flux produced by the gap commutator is
\begin{equation}
  \partial_E\!\left[
    \frac{\Delta}{E}\dot\Delta\,\DOS\,\fL
  \right]
  =
  \partial_E\!\left[\dot\Delta\,N_2\,\fL\right],
  \label{eq:missing_flux_N2}
\end{equation}
where \(N_2=(\Delta/E)\DOS\) on the positive-energy BCS branch. The
branch-Boltzmann route in \cref{sec:branch_alternative} produces this same
term before angular averaging, so the two routes now agree; this
subsection verifies that no other piece of the Wigner/Moyal reduction
double-counts or modifies it at leading order.

\paragraph{Higher Moyal orders in the gap commutator.} The all-orders
Moyal product for a gap operator
\(A_\Delta=-\hatD=\ii\Delta(t)\tau_2\), which has no \(E\) dependence, is
\begin{equation}
  A_\Delta\star B
  =
  \sum_{n=0}^{\infty}
  \frac{1}{n!}\left(\frac{\ii\hbar}{2}\right)^n
  (-1)^n
  \bigl(\partial_t^n A_\Delta\bigr)\,
  \partial_E^n B,
  \qquad
  B\star A_\Delta
  =
  \sum_{n=0}^{\infty}
  \frac{1}{n!}\left(\frac{\ii\hbar}{2}\right)^n
  \bigl(\partial_E^n B\bigr)\,
  \bigl(\partial_t^n A_\Delta\bigr),
  \label{eq:gap_moyal_all_orders}
\end{equation}
with the displayed left--right matrix ordering. Therefore
\begin{equation}
  [A_\Delta,B]_\star
  =
  \sum_{n=0}^{\infty}
  \frac{1}{n!}\left(\frac{\ii\hbar}{2}\right)^n
  \left[
    (-1)^n\bigl(\partial_t^n A_\Delta\bigr)\partial_E^n B
    -
    \bigl(\partial_E^n B\bigr)\bigl(\partial_t^n A_\Delta\bigr)
  \right].
  \label{eq:gap_comm_all_orders}
\end{equation}
For the longitudinal, charge-balanced symbol used above,
\(B=\hatg_0^K=(2\DOS\tau_3+2\ii N_2\tau_2)\fL\), every
\(\partial_E^n B\) carries a \(\tau_3\) piece and a \(\tau_2\) piece. Even
\(n\) terms contain the commutators \([\tau_2,\tau_3]\propto\tau_1\)
(traceless) and \([\tau_2,\tau_2]=0\), so they contribute no plain-trace
scalar at any even order. Odd \(n\) terms contain the anticommutators
\(\{\tau_2,\tau_3\}=0\) and \(\{\tau_2,\tau_2\}=2\,\openone\): the
\(\tau_2\) piece of \(\partial_E^n B\) survives the trace. At \(n=1\) this
is exactly the flux \cref{eq:trace_gap}; at higher odd \(n\) the
contribution carries \(\hbar^n\Delta^{(n)}\)
(\(\ddot\Delta,\Delta^{(3)},\ldots\)), genuinely higher order in the
adiabatic expansion. The leading \(O(\hbar\dot\Delta)\) flux is therefore
produced once, at \(n=1\), with no double counting from the Moyal tail.

The \(E\tau_3\) part gives no hidden higher-order source either. Since
\(\partial_E^2(E\tau_3)=0\) and \(\partial_t(E\tau_3)=0\), its Moyal series
terminates exactly at
\begin{equation}
  \ii[E\tau_3,B]_\star
  =
  \ii E[\tau_3,B]
  -
  \frac{\hbar}{2}\{\tau_3,\partial_t B\},
  \label{eq:E_moyal_exact_test}
\end{equation}
which is the time-derivative term already evaluated in
\cref{eq:trace_time_deriv}.

\paragraph{Time-Moyal corrections to the matrix current.} The only place in
the matrix current where a time-dependent spectrum could have been hidden is
the spectral-anomalous drift line of \cref{eq:matrix_current_general},
\[
  \hatg^{R}\widetilde\bnabla\hatg^{R}\cdot\hat h
  -
  \hat h\cdot\hatg^{A}\widetilde\bnabla\hatg^{A},
\]
with all products promoted to \(\star\)-products. For spatially uniform
\(\Delta(t)\), however, \(\widetilde\bnabla\hatg^{R/A}=0\) as a Wigner symbol,
and every \(E\)- or \(t\)-derivative of this zero symbol is also zero. The
time-Moyal product of the drift term therefore remains zero. Moyal
corrections to the first-line diffusive flux can only modify the spatial
current, producing terms inside
\(\widetilde\bnabla\!\cdot(\cdots)\) that contain spatial gradients of
\(\fL\) or its \(E,t\) derivatives. They cannot produce the homogeneous
energy-space divergence in \cref{eq:missing_flux_N2}, which is present even
when \(\widetilde\bnabla\fL=0\).

\paragraph{First Wigner correction to the Keldysh ansatz and projection.}
One might also ask whether replacing the \(\star\)-products in
\(\hatg^K=\hatg^R\star\hat h-\hat h\star\hatg^A\) by ordinary products before
projecting shifts the result. For \(\hat h=\fL\openone\), the first correction is
\begin{equation}
  \delta_\star\hatg^K
  =
  \frac{\ii\hbar}{2}
  \left[
    (\partial_E\hatg^R+\partial_E\hatg^A)\,\partial_t\fL
    -
    (\partial_t\hatg^R+\partial_t\hatg^A)\,\partial_E\fL
  \right].
  \label{eq:gK_star_correction_general}
\end{equation}
Using the BCS matrices in \cref{eq:gRA_BCS},
\(\hatg^R+\hatg^A=0\) above the gap, so
\begin{equation}
  \delta_\star\hatg^K
  =
  0 :
  \label{eq:gK_star_correction_BCS}
\end{equation}
the first Wigner correction to the distribution ansatz vanishes
identically for the charge-balanced symbol. (Its plain-trace contribution
through the ordinary commutator \(\ii[\hat L_0,\delta_\star\hatg^K]\)
would vanish by cyclicity in any case,
\(\Tr[\ii[\hat L_0,\delta_\star\hatg^K]]=0\).) The ansatz correction
therefore neither supplies nor disturbs the leading
\(O(\hbar\dot\Delta)\) energy-space flux.

There is also no separate Moyal correction to the projector itself for the
matrix kinetic equations used here: the projectors are the constant Nambu
matrices \(\openone\) and \(\tau_3\), so
\(\openone\star X=\openone X\) and \(\tau_3\star X=\tau_3X\) exactly. An
energy- and \(\Delta\)-dependent Bogoliubov branch projector would have its
own Wigner corrections and does expose the branch kinematics
\(\dot E=(\Delta/E)\dot\Delta\), but inserting that projector changes the
order of operations to the branch-Boltzmann route of
\cref{sec:branch_alternative}. It is not a completion of the plain-trace
matrix-Usadel projection.

\paragraph{Why the drive survives the projection.} The structural reason
the flux appears becomes visible once
the spectral commutator is expanded through the ansatz
\(\hatg_0^K=\hatg^R\star\hat h-\hat h\star\hatg^A\) rather than from the
instantaneous symbol. Isolating the Moyal commutator of the spectral operator
with the distribution matrix, for \(\hat h=\fL\openone\),
\begin{equation}
  [\hat L_0,\hat h]_\star
  =
  \ii\hbar\,\partial_t\fL\;\tau_3
  +
  \hbar\,\dot\Delta\,\partial_E\fL\;\tau_2
  +\cdots,
  \label{eq:drive_L0h}
\end{equation}
with \(\hat L_0=E\tau_3+\ii\Delta\tau_2\) the spectral operator in the gap
gauge of \cref{sec:conventions}; the omitted terms are higher odd Moyal
orders, all \(\propto\tau_2\) (the even orders cancel, as the matrix
\(\star\)-commutator with \(\hat h\propto\openone\) reduces to the scalar
Moyal bracket). The \(\tau_2\) term \emph{is} the spectral-flow drive.

The longitudinal Keldysh projection retains it. Expanding
\(\ii[\hat L_0,\hatg_0^K]_\star\) with the Leibniz rule for the
\(\star\)-commutator, the term that carries the distribution-matrix drive
\([\hat L_0,\hat h]_\star\) appears sandwiched between the spectral
propagators as
\(\hatg^R\star[\hat L_0,\hat h]_\star-[\hat L_0,\hat h]_\star\star\hatg^A\),
and to leading order its \(\tau_2\) part carries the matrix structure
\begin{equation}
  \hatg^R\tau_2-\tau_2\hatg^A=2\ii N_2\,\openone\;\neq\;0,
  \label{eq:drive_cancels}
\end{equation}
because the corrected \(\hatg^A=-\hatg^R\) makes the anomalous components
of \(\hatg^R\) and \(\hatg^A\) \emph{opposite}, not shared
[\cref{eq:gRA_BCS}]. The drive's sandwich weight is proportional to the
identity, so it survives the plain trace, with scalar weight
\(2N_2\)---precisely assembling the flux
\(\partial_E[\dot\Delta N_2\fL]\) of \cref{eq:trace_gap}. The \(\tau_3\)
time-derivative drive likewise survives,
\(\hatg^R\tau_3-\tau_3\hatg^A=2\DOS\,\openone\ne0\), giving
the \(\partial_t(\DOS\fL)\) term of \cref{eq:trace_time_deriv}.

(For the opposite advanced-propagator sign the
anomalous components are shared rather than opposite and the weight
\cref{eq:drive_cancels} vanishes, annihilating the drive---another
statement of the requirements that fix main-text~\cref{M-eq:bcs_advanced_matrix}.
The fixed-\(E\) plain-trace projection generates the energy-space flux
\cref{eq:missing_flux_N2} directly, and the moving-coordinate
constructions below are its kinematic re-expressions.)
The constant-trace matrix reduction is therefore complete as a fixed-\(E\)
conservation law; the next subsections rewrite that law on the moving
quasiparticle energy coordinate.

\subsubsection{Coordinate-lift attempt: the moving quasiparticle energy shell}
\label{sec:coordinate_lift}

The negative result above is a statement about a specific projection: take the
Keldysh--Usadel equation written in the Wigner energy \(E\), project with the
constant longitudinal Nambu trace, and look for a local matrix term
proportional to \(\dot\Delta\,\partial_E\fL\). Within that operation the term
is not present. There is nevertheless a precise kinematic way to recover the
canonical spectral-flow flux once the static Usadel continuity law is lifted
from a fixed spectrum to an instantaneous quasiparticle spectrum.

For a spatially uniform but time-dependent BCS gap, label the positive-energy
states by the normal-state dispersion magnitude
\begin{equation}
  \xi \equiv \sqrt{E^2-\Delta^2(t)} .
  \label{eq:xi_label}
\end{equation}
The branch energy at fixed \(\xi\) is
\begin{equation}
  E_\xi(t)=\sqrt{\xi^2+\Delta^2(t)},
  \qquad
  \left(\partial_t E_\xi\right)_\xi
  =
  \frac{\Delta}{E}\dot\Delta
  \equiv u_E(E,t).
  \label{eq:moving_energy_velocity}
\end{equation}
The BCS density of states is exactly the Jacobian between the invariant label
and the quasiparticle energy,
\begin{equation}
  \DOS(E,\Delta)=\frac{\partial\xi}{\partial E}
  =\frac{E}{\sqrt{E^2-\Delta^2}},
  \qquad
  \DOS\,\dd E=\dd\xi ,
  \label{eq:DOS_jacobian}
\end{equation}
up to the conventional factor counting the two \(\pm\xi\) branches, which is
absorbed into the normal-state density of states. Consequently a branch
occupation \(F(\xi,\br,t)\equiv\fL(E_\xi(t),\br,t)\) obeys
\begin{equation}
  \left(\partial_t F\right)_\xi
  =
  \left(\partial_t+u_E\,\partial_E\right)\fL .
  \label{eq:xi_time_derivative}
\end{equation}
The same statement can be written as a conservative law for the spectral
density \(n_E=\DOS\fL\). Using
\begin{equation}
  \partial_t\DOS+\partial_E(u_E\DOS)=0,
  \label{eq:DOS_continuity_from_jacobian}
\end{equation}
which follows directly by differentiating the Jacobian
\(\DOS=\partial\xi/\partial E\) at fixed \(\xi\). In the matrix language the
same identity is the adiabatic spectral equation: to first Moyal order,
\begin{equation}
  [\,\hat L_0,\hatg^R\,]_\star
  =
  \ii\hbar\,
  \bigl[\partial_t\DOS+\dot\Delta\,\partial_E N_2\bigr]\openone
  =
  0,
  \label{eq:gR_star_DOS_continuity}
\end{equation}
because \(u_E\DOS=\dot\Delta N_2\). Thus the coordinate lift is not an
external phenomenological addition; it is the retarded spectral flow of the
instantaneous BCS propagator written in the moving energy coordinate. Using
\cref{eq:DOS_continuity_from_jacobian}, one finds
\begin{equation}
  \partial_t(\DOS\fL)
  +\partial_E(u_E\DOS\fL)
  =
  \DOS\left(\partial_t+u_E\partial_E\right)\fL .
  \label{eq:coordinate_lift_identity}
\end{equation}
This identity is the same energy-space flux that the fixed-\(E\)
plain-trace projection produces directly (\cref{eq:trace_gap}); here it
arises kinematically, from asking
that the scalar density proven by the static trace, \(\DOS\fL\), be transported
between energy bins when the map \(\xi\leftrightarrow E\) changes in time.

Written with the (undressed) dirty-limit Usadel current, the equation reads
\begin{equation}
  \boxed{\,
  \hbar\left[
    \partial_t(\DOS\fL)
    +\partial_E\!\left(
      \frac{\Delta}{E}\dot\Delta\,\DOS\fL
    \right)
  \right]
  =
  \hbar\DN\,\widetilde\bnabla^{2}\fL
  +\hbar J_1[f,n].
  \,}
  \label{eq:coordinate_lifted_A1}
\end{equation}
Equivalently,
\begin{equation}
  \hbar\,\DOS
  \left[
    \partial_t
    +\frac{\Delta}{E}\dot\Delta\,\partial_E
  \right]\fL
  =
  \hbar\DN\,\widetilde\bnabla^{2}\fL
  +\hbar J_1[f,n].
  \label{eq:coordinate_lifted_A1_advective}
\end{equation}
After \(\fL=1-2f\), this is the target time-dependent A1 scalar equation.

The construction shows that the spectral-flow term is not an arbitrary
extra term: it is the Liouville/Jacobian correction required when the
conserved scalar selected by the dirty matrix projection, \(\DOS\fL\), is
expressed in the moving local quasiparticle energy coordinate---fixed
energy bins are not fixed quasiparticle states when \(\Delta(t)\)
changes. Since the fixed-\(E\) projection produces the same conservative
form directly (\cref{sec:trace_projection}), the lift is its kinematic
interpretation, not an independent completion. The weak-form version in
\cref{sec:moving_bin_projection} builds the statement into the scalar
projection itself, and the branch-projector form in
\cref{sec:branch_projector_lift} expresses it in an instantaneous BdG
basis. This is a question of representation, not of missing physics: a
moving-\(\xi\) or gauge-invariant-energy Wigner representation would make the
spectral flow manifest before the constant-trace reduction, but it yields no
term the fixed-\(E\) reduction lacks. (Symbolic verification confirms the companion statement for
the spatially varying gap---the fixed-\(E\) longitudinal matrix current is
the undressed \(\bnabla\fL\)
with no spurious \(\bnabla\Delta\,\partial_E\) drift and no
longitudinal--transverse mixing.) The calculation in
\cref{sec:tdep_candidate_tests} gives the same content from the
same fixed-\(E\) equation, as expected.

\subsubsection{Weak-form projection on moving spectral bins}
\label{sec:moving_bin_projection}

The coordinate lift can be made less ad hoc by formulating the scalar
projection in weak form. Choose a fixed interval
\(\mathcal I_\xi=[\xi_a,\xi_b]\) of the invariant normal-state label, and let
\(E_a(t)=E_{\xi_a}(t)\), \(E_b(t)=E_{\xi_b}(t)\). The longitudinal content of
that interval is
\begin{equation}
  Q_{\mathcal I}(\br,t)
  =
  \int_{\xi_a}^{\xi_b}\dd\xi\,
    F(\xi,\br,t)
  =
  \int_{E_a(t)}^{E_b(t)}\dd E\,
    \DOS(E,t)\,\fL(E,\br,t),
  \label{eq:moving_bin_charge}
\end{equation}
where \(F(\xi,\br,t)=\fL(E_\xi(t),\br,t)\). Differentiating at fixed
\(\xi_a,\xi_b\) gives the Reynolds-transport form
\begin{equation}
  \partial_t Q_{\mathcal I}
  =
  \int_{E_a}^{E_b}\dd E\,
  \left[
    \partial_t(\DOS\fL)
    +\partial_E(u_E\DOS\fL)
  \right],
  \qquad
  u_E=\frac{\Delta}{E}\dot\Delta .
  \label{eq:moving_bin_reynolds}
\end{equation}
Thus the energy-space flux is the boundary motion of a spectral
control volume whose endpoints are fixed quasiparticle states.

The dirty matrix reduction supplies the spatial part of the weak law: after
elastic impurity isotropization, the longitudinal current carried by each
instantaneous energy shell is
\begin{equation}
  \mathbf j_L(E,\br,t)
  =
  -\DN\,\widetilde\bnabla\fL(E,\br,t),
  \label{eq:moving_bin_jL}
\end{equation}
the same undressed current derived in \cref{sec:DL_DT}. Assuming the spectrum changes
adiabatically on the elastic isotropization time, so that this instantaneous
current remains the dirty-limit spatial current inside each moving shell, the
weak conservation statement is
\begin{equation}
  \hbar\,\partial_t Q_{\mathcal I}
  =
  \int_{E_a}^{E_b}\dd E\,
  \left\{
    \hbar\DN\,\widetilde\bnabla^{2}\fL
    +\hbar J_1[f,n]
  \right\}.
  \label{eq:moving_bin_weak_law}
\end{equation}
Because the interval \(\mathcal I_\xi\) is arbitrary, localizing
\cref{eq:moving_bin_weak_law} in \(E\) gives exactly
\cref{eq:coordinate_lifted_A1}. In this formulation the spectral-flow
term appears as the boundary term required by the physical choice of control
volume---bins of fixed quasiparticle state label, not bins of fixed Wigner
energy---in exact agreement with the pointwise fixed-\(E\) trace equation
\cref{eq:expected_conservative}.

This also states the cost of the step. The argument is intrinsic to the
scalar, energy-shell projection, but it still uses the instantaneous BCS
spectral map and the static dirty-limit current. It assumes no nonadiabatic
spectral corrections, no Landau-Zener-like branch conversion, and no unresolved
charge-imbalance dynamics beyond the \(J_1\) collision bookkeeping. For a
spatially varying gap, fixed-\(\xi\) and fixed-\(E\) spatial gradients also
differ; the cancellation of the corresponding \(\nabla\Delta\,\partial_E\)
terms must be checked together with the full matrix current and boundary
conditions. Those extensions are outside the uniform-gap time-flow closure
established here.

\subsubsection{Adiabatic branch-projector form}
\label{sec:branch_projector_lift}

The moving-bin argument can also be phrased as a projection onto instantaneous
BdG branches. This is the closest version, within this appendix, to building the
energy-shell motion into the projection itself.

Work temporarily in a Hermitian gap gauge, gauge-equivalent to the Nambu basis
used above, with local BdG Hamiltonian
\begin{equation}
  H_\xi(t)=\xi\,\tau_3+\Delta(t)\,\tau_1,
  \qquad
  E_\xi(t)=\sqrt{\xi^2+\Delta^2(t)} .
  \label{eq:bdg_branch_hamiltonian}
\end{equation}
The instantaneous positive/negative branch projectors are
\begin{equation}
  P_\pm(\xi,t)
  =
  \frac{1}{2}
  \left[
    \openone
    \pm
    \frac{H_\xi(t)}{E_\xi(t)}
  \right],
  \qquad
  P_\pm^2=P_\pm,\quad P_+P_-=0 .
  \label{eq:bdg_branch_projectors}
\end{equation}
For a charge-balanced longitudinal distribution, the branch-diagonal
occupation extracted from the distribution matrix is
\begin{equation}
  F_+(\xi,\br,t)
  =
  \Tr\!\left[P_+\,\hat h_\xi\,P_+\right]
  =
  \fL(E_\xi(t),\br,t)
  \qquad(\fT=0),
  \label{eq:branch_projected_fL}
\end{equation}
up to the trivial normalization of the rank-one projector. In the
instantaneous eigenbasis the projected kinetic equation contains a covariant
time derivative,
\begin{equation}
  D_t \hat h_\xi
  =
  \partial_t \hat h_\xi
  +[\mathcal A_t,\hat h_\xi],
  \qquad
  \mathcal A_t=U^\dagger\partial_t U ,
  \label{eq:branch_covariant_derivative}
\end{equation}
where \(U(\xi,t)\) diagonalizes \(H_\xi(t)\). The connection
\(\mathcal A_t\) is off-diagonal between the \(+\) and \(-\) branches for a
real \(s\)-wave gap. Its commutator does not change the diagonal longitudinal
occupation when \(\hat h_\xi\propto\openone\) in branch space, and the
remaining off-diagonal branch coherence is higher order in the adiabatic
parameter \(\hbar\dot\Delta/E_\xi^2\). To leading adiabatic order,
\begin{equation}
  (D_t \hat h_\xi)_{++}
  =
  \left(\partial_t F_+\right)_\xi
  =
  \left[
    \partial_t
    +\left(\partial_t E_\xi\right)_\xi\partial_E
  \right]\fL
  =
  \left[
    \partial_t
    +\frac{\Delta}{E}\dot\Delta\,\partial_E
  \right]\fL .
  \label{eq:branch_projector_spectral_flow}
\end{equation}
Thus an energy- and gap-dependent branch projector exposes, in advective
form, the same spectral-flow derivative that the constant Nambu trace
produces in conservative form (\cref{eq:expected_conservative}). The
operation differs---\(P_+(\xi,t)\) is not the constant longitudinal
projector used in \cref{sec:trace_projection}; it carries the time
dependence of the local BdG eigenvectors and the moving eigenvalue
\(E_\xi(t)\)---but the content is the same.

Combining this branch-projected time derivative with the dirty matrix-Usadel
spatial current from \cref{sec:DL_DT} gives
\begin{equation}
  \hbar\,\DOS
  \left[
    \partial_t
    +\frac{\Delta}{E}\dot\Delta\,\partial_E
  \right]\fL
  =
  \hbar\DN\,\widetilde\bnabla^{2}\fL
  +\hbar J_1[f,n],
  \label{eq:branch_projector_A1}
\end{equation}
which is identical to \cref{eq:coordinate_lifted_A1_advective}. This is a
hybrid but controlled construction: the branch projector fixes the moving
spectral time coordinate, while the spatial diffusion coefficient and the
undressed longitudinal flux remain those of the dirty matrix-Usadel reduction.
It is therefore different from the clean branch-Boltzmann route in
\cref{sec:branch_alternative}, where the spatial mobility is also projected
before impurity isotropization.

This settles the formal status of the construction. A representation-intrinsic
derivation in a moving-\(\xi\), gauge-invariant-energy, or branch-covariant
Wigner frame would reproduce \cref{eq:branch_projector_A1} with the spectral
flow manifest before the scalar limit---but it is an exact re-expression of the
fixed-\(E\) result, not a route to new terms. The fixed-\(E\) constant-trace
reduction already yields the complete content (the branch-diagonal
charge-balanced sector reduces to \cref{eq:branch_projector_A1}; the off-diagonal
coherences are suppressed by the adiabatic parameter of
\cref{eq:branch_coherence_estimate}); the moving-frame version differs from it
only by the exact bin relabeling. What lies beyond the present assumptions is
no longer fully open: the supercurrent-induced longitudinal--transverse mixing,
the non-BCS proximity dressings, and the leading nonadiabatic corrections are
established in \cref{app:supercurrent,app:proximity,app:nonadiabatic};
genuinely open beyond those bounds is the collision-level physics
(nonadiabatic memory, drive sidebands, a gap changing appreciably during a
collision) and self-consistent gap dynamics.

\paragraph{Size of the discarded branch coherence.} The adiabatic condition
can be made explicit. Let
\begin{equation}
  \cos\theta_\xi=\frac{\xi}{E_\xi},
  \qquad
  \sin\theta_\xi=\frac{\Delta}{E_\xi},
  \qquad
  U_\xi=\exp\!\left(-\frac{\ii}{2}\theta_\xi\tau_2\right),
  \label{eq:branch_rotation}
\end{equation}
so that \(U_\xi^\dagger H_\xi U_\xi=E_\xi\tau_3\). Then
\begin{equation}
  \mathcal A_t
  =
  U_\xi^\dagger\partial_t U_\xi
  =
  -\frac{\ii}{2}\dot\theta_\xi\tau_2,
  \qquad
  \dot\theta_\xi
  =
  \left(\partial_t\theta_\xi\right)_\xi
  =
  \frac{\xi\,\dot\Delta}{E_\xi^2}.
  \label{eq:branch_connection_size}
\end{equation}
In the branch eigenbasis, an off-diagonal distribution component
\(h_{+-}\) obeys schematically
\begin{equation}
  \left(
    \partial_t+\Gamma_{+-}+\frac{2\ii E_\xi}{\hbar}
  \right)h_{+-}
  =
  -(\mathcal A_t)_{+-}\,(h_{--}-h_{++})
  +\cdots ,
  \label{eq:branch_coherence_equation}
\end{equation}
where \(\Gamma_{+-}\) represents any dephasing from impurity, inelastic, or
order-parameter relaxation physics not kept explicitly in the scalar sector,
and the omitted terms contain spatial gradients and collisions. In the
adiabatic regime \(|\partial_t h_{+-}|\ll 2E_\xi|h_{+-}|/\hbar\),
\begin{equation}
  h_{+-}
  \sim
  \frac{\ii\hbar\,(\mathcal A_t)_{+-}}{2E_\xi-\ii\hbar\Gamma_{+-}}\,
  (h_{--}-h_{++})
  =
  O\!\left(
    \frac{\hbar|\xi\dot\Delta|}{E_\xi^3}
  \right)
  (h_{--}-h_{++}).
  \label{eq:branch_coherence_estimate}
\end{equation}
Thus the branch-projector reduction is controlled when
\(\hbar|\xi\dot\Delta|/E_\xi^3\ll1\), and it is even safer in the
charge-balanced longitudinal sector where the branch-antisymmetric source
\(h_{--}-h_{++}\) is absent or rapidly relaxed into the transverse channel.
This estimate is the branch-basis version of the assumption already implicit
in using instantaneous BCS spectral functions: the gap changes slowly compared
with the local quasiparticle splitting \(2E_\xi/\hbar\).

\subsubsection{Conservative-form longitudinal kinetic equation:
\texorpdfstring{\(p=1\)}{p=1}}
\label{sec:conservative_form}

Combining \cref{eq:trace_diffusion}, \cref{eq:trace_time_deriv},
\cref{eq:trace_gap}, and \cref{eq:J1_def} with the K-block kinetic
equation \cref{eq:K_block_kinetic} (after taking plain trace and
dividing by \(4\)) gives the longitudinal kinetic equation that the
plain-trace projection produces:
\begin{equation}
  \boxed{\,
  \hbar\left[
  \partial_t\bigl(\DOS\,\fL\bigr)
  +
  \partial_E\!\left(\frac{\Delta}{E}\dot\Delta\,\DOS\,\fL\right)
  \right]
  =
  \hbar\DN\,\widetilde\bnabla^{2}\fL
  + \hbar J_1[f,n].
  \,}
  \label{eq:full_kinetic_conservative}
\end{equation}
This equation has the canonical structure --- the conserved spectral density
on the LHS is \(\DOS\fL\), transported by the spectral-flow energy flux,
the spatial flux is undressed (\(\mathcal D_L=1\) above the gap), and the
collision integral \(J_1\) is the standard inelastic kernel. In the
\(L_{p,q}\) classification of the main text, the spatial operator on
the RHS is \(L_{1,0}\), i.e., \((p,q)=(1,0)\): time-weight exponent
\(p=1\), flux-dressing exponent \(q=0\). For static \(\Delta\),
\(\partial_t\DOS=0\), the energy flux drops, and
\cref{eq:full_kinetic_conservative} collapses to
\begin{equation}
  \boxed{\,
  \hbar\,\DOS(E,\br)\,\partial_t\fL
  =
  \hbar\DN\,\widetilde\bnabla^{2}\fL
  + \hbar J_1[f,n]
  \quad(\text{static }\Delta).
  \,}
  \label{eq:static_kinetic_fL}
\end{equation}

For \emph{time-dependent} \(\Delta\),
\cref{eq:full_kinetic_conservative} is already the canonical
scalar-Boltzmann statement: the BCS DOS continuity
\cref{eq:DOS_continuity} rewrites its LHS as the advective derivative
\(\DOS[\partial_t+(\Delta/E)\dot\Delta\,\partial_E]\fL\). The coordinate
lift of \cref{sec:coordinate_lift}, the moving-bin weak projection of
\cref{sec:moving_bin_projection}, and the branch-projector form of
\cref{sec:branch_projector_lift} reproduce the same equation as the
Liouville statement for the moving quasiparticle energy coordinate; they
are kinematic re-expressions of
\cref{eq:full_kinetic_conservative}, not additional ingredients.

\subsubsection{The transverse kinetic equation}
\label{sec:transverse_projection}

Before setting \(\fT\to0\), we assemble the transverse (charge-imbalance)
kinetic equation produced by the same matrix machinery. Three ingredients
enter: the dressed transverse flux already computed, the first Wigner
correction to the Keldysh ansatz---nonzero in the \(\fT\) sector, unlike
the longitudinal one---and the channel projection that converts the matrix
residual into a scalar equation. The result,
\cref{eq:fT_kinetic_conservative}, is the charge-mode companion of
\cref{eq:full_kinetic_conservative}: the same conservative spectral-flow
pair, plus a single additional coherent conversion term with a kinematic
meaning, and no \(\fL\) source.

The spatial flux follows from \cref{eq:matrix_current_BCS}:
\begin{equation}
  \Tr\!\left[
    \tau_3
    \bigl(\checkg_0\widetilde\bnabla\checkg_0\bigr)^K
  \right]
  =
  4\,\DOS^2\,\widetilde\bnabla\fT ,
  \label{eq:traceT_diffusion}
\end{equation}
so the physical charge current is dressed,
\(\mathbf j_T=-\DN\DOS^2\widetilde\bnabla\fT\), as already found in
\cref{eq:jT_BCS}. The \(E\tau_3\) Moyal correction contributes the time
term
\begin{equation}
  -\frac{\hbar}{2}
  \Tr\!\left[
    \tau_3\{\tau_3,\partial_t\hatg_0^K\}
  \right]
  =
  -\hbar\,\Tr[\partial_t\hatg_0^K]
  =
  -4\hbar\,\partial_t(\DOS\fT),
  \label{eq:traceT_time_deriv}
\end{equation}
using \(\Tr[\hatg_0^K]=4\DOS\fT\). In the longitudinal sector the
analogous pair of trace statements, together with the gap-drive term,
already assembles into the kinetic equation, because the first Wigner
correction to the ansatz
\(\hatg_0^K=\hatg^R\circ\hat h-\hat h\circ\hatg^A\) vanishes identically
in the \(\fL\) sector. In the transverse sector it does not:
\begin{equation}
  \delta_\circ\hatg_0^K\big|_{\fT}
  =
  \ii\hbar\,\frac{E}{W^3}
  \bigl(\Delta\,\partial_t\fT+E\dot\Delta\,\partial_E\fT\bigr)\,\tau_1 ,
  \qquad W=\sqrt{E^2-\Delta^2},
  \label{eq:deltagK_fT}
\end{equation}
a pure \(\tau_1\) (branch-off-diagonal) piece, nonzero whenever \(\fT\) or
the gap is time dependent. Fed through \([\hat L_0,\tau_1]\), it adds
\(\fT\)-sector terms to the \(\tau_3\)-trace of the kinetic commutator
beyond \cref{eq:traceT_time_deriv}, so
\cref{eq:traceT_diffusion,eq:traceT_time_deriv} must not be combined at
face value into a transverse kinetic equation; the consistent assembly is
the channel analysis next.

\paragraph{Kernel channels and the scalar projection.}
Above the gap, with a real order parameter, the map
\(X\mapsto[\hat L_0,X]\) splits Nambu space into a two-dimensional kernel
and a two-dimensional image,
\begin{equation}
  \ker=\mathrm{span}\{\openone,\hatg^0\},
  \quad
  \operatorname{im}=\mathrm{span}\{\tau_1,\hat m\},
  \quad
  \hat m\equiv N_2\tau_3+\ii\DOS\tau_2,
  \label{eq:kernel_channels}
\end{equation}
where \(\hatg^0=\DOS\tau_3+\ii N_2\tau_2\) is the normalized spectral
matrix and \([\hat L_0,\tau_1]=2W\hat m\), \([\hat L_0,\hat m]=2W\tau_1\).
Under the bilinear trace pairing \(\langle X,Y\rangle=\tfrac12\Tr[XY]\)
the kernel and image are orthogonal, with
\(\tau_3=\DOS\,\hatg^0-N_2\,\hat m\) and
\(\tau_2=\ii N_2\,\hatg^0-\ii\DOS\,\hat m\). The symbolic verification
establishes that, to first Moyal order and for a spatially uniform gap,
the full kinetic commutator lies \emph{exactly} in the kernel:
\begin{equation}
\begin{aligned}
  \ii\bigl[\hat L_0,\hatg_0^{K}\bigr]_{\circ}
  &=
  A\,\openone+B\,\hatg^0,
  \\
  A&=-2\hbar\bigl[\partial_t(\DOS\fL)+\partial_E(\dot\Delta N_2\fL)\bigr],
  \\
  B&=-2\hbar\,\DOS^2
  \left[\partial_t
  +\frac{\Delta\dot\Delta}{E}\,\partial_E
  +\frac{\Delta\dot\Delta}{E^{2}}\right]\fT ,
\end{aligned}
  \label{eq:kernel_residual}
\end{equation}
with identically vanishing \(\tau_1\) and \(\hat m\) components (the
spatially varying case is taken up below). The scalar kinetic equations
are the kernel-channel coefficients: the identity channel
\(A=\tfrac12\Tr[\,\cdot\,]\) is the longitudinal projection, already
\cref{eq:full_kinetic_conservative}, and the spectral channel
\(B=\tfrac12\Tr[\hatg^0\,\cdot\,]\) is the transverse one. Because the
factorization
\([\hat L_0,\hatg^K]_\star
=\hatg^R\star[\hat L_0,\hat h]_\star-[\hat L_0,\hat h]_\star\star\hatg^A\)
holds to this order (\cref{eq:specR,eq:specA}), these kernel channels are
the two real scalar equations carried by the distribution-matrix form of
the kinetic theory (Belzig et~al.~\cite{belzig1999}, their Eq.~30); the
image channels carry slaved spectral content, and no kinetic equation
lives there. The \(\tau_3\)-trace of \cref{eq:kernel_residual} equals
\(2\DOS B\): an \(\DOS\)-weighted copy of the transverse kernel channel.
Projection weights of this kind---and the swap whereby the \emph{plain}
trace of \(\hatg_0^K\) carries the transverse density \(4\DOS\fT\) while
the plain trace of the \emph{equation} is longitudinal---are why pairing
\cref{eq:traceT_time_deriv} with \cref{eq:traceT_diffusion} at face value
would miscount the transverse time weight by \(\DOS^{2}\).

\paragraph{The transverse kinetic equation.}
Adding the gradient flux and the collision integral projected onto the
same channel,
\begin{equation}
  4\hbar\,\DOS\,J_T[\fL,\fT,n](E,\br)
  \equiv
  \ii\,\Tr\!\bigl[\hatg^0\,\Icoll[\checkg_0]\bigr],
  \label{eq:JT_def}
\end{equation}
the transverse companion of \cref{eq:full_kinetic_conservative} is
\begin{equation}
  \boxed{\,
  \hbar\left[
  \partial_t\bigl(\DOS\,\fT\bigr)
  +\partial_E\!\left(\frac{\Delta}{E}\dot\Delta\,\DOS\,\fT\right)
  +\frac{\Delta\dot\Delta}{E^{2}}\,\DOS\,\fT
  \right]
  =
  \hbar\DN\,\widetilde\bnabla^{2}\fT
  +\hbar\,J_T
  \quad(E>\Delta),
  \,}
  \label{eq:fT_kinetic_conservative}
\end{equation}
for a spatially uniform gap (the spatially varying flux is given below).
For static \(\Delta\) it collapses to
\begin{equation}
  \hbar\,\DOS\,\partial_t\fT
  =
  \hbar\DN\,\widetilde\bnabla^{2}\fT
  +\hbar\,J_T
  \qquad(\text{static }\Delta),
  \label{eq:fT_static}
\end{equation}
with the same time weight \(p=1\) and the same undressed uniform-bulk
flux as the longitudinal pair \cref{eq:static_kinetic_fL}: in uniform
bulk both modes diffuse at the rate \(\DN/\DOS(E)\). The \(\DOS^2\)
transverse dressing of \cref{eq:traceT_diffusion} is the dressing of the
physical charge current \(\mathbf j_T\), not of the kinetic-equation flux
in this normalization. Accordingly \(\partial_t(\DOS\fT)\) and
\(\widetilde\bnabla\!\cdot\!\mathbf j_T\) do not balance by themselves:
the difference is coherent charge exchange with the condensate---the
quasiparticle charge alone is not locally conserved, and the balance is
closed by the supercurrent response.

\paragraph{Kinematic reading and particle--hole symmetry.}
The structure new relative to the longitudinal sector is the conversion
term \((\Delta\dot\Delta/E^{2})\DOS\fT\). Multiplying
\cref{eq:fT_kinetic_conservative} by \(E\) shows that \(u\equiv E\fT\)
obeys exactly the longitudinal operator
\(\partial_t(\DOS u)+\partial_E((\Delta/E)\dot\Delta\,\DOS u)\), so the
collisionless solutions are
\begin{equation}
  \fT(E,t)=\frac{F(\xi)}{E},
  \qquad
  \xi=\sqrt{E^{2}-\Delta(t)^{2}},
  \label{eq:fT_frozen}
\end{equation}
for arbitrary \(F\): the charge-mode analogue of the frozen-shell
longitudinal solutions \(\fL=G(\xi)\). The kinematics is the adiabatic
theorem plus charge dilution. Per-shell branch-odd occupations are frozen
on the moving shells \(E=\sqrt{\xi^2+\Delta(t)^2}\) exactly as in the
energy mode, while the quasiparticle effective charge \(q=\xi/E\) is
diluted as the gap moves; with the charge density normalized as
\(\int\DOS\fT\,\dd E\), the mode amplitude falls as \(1/E\) along the
shell, the difference flowing coherently to the condensate at the
collisionless rate \(-\dd\ln q/\dd t=\Delta\dot\Delta/E^{2}\). Two
structural negatives are explicit in \cref{eq:kernel_residual}: there is
no \(\bigl[2\Delta^2/\sqrt{E^2-\Delta^2}\bigr]\fT\)-type coherent
relaxation above the gap, and there is no \(\fL\) source anywhere in the
transverse channel---a spatially uniform real \(\dot\Delta\)
redistributes existing charge imbalance but cannot create it, as
particle--hole symmetry requires. Both statements hold with \(\fT\)
retained at first Moyal order, not merely at \(\fT=0\).

\paragraph{Relaxation under depairing, and supercurrent coupling.}
Belzig et~al.'s relaxation coefficient (their Eq.~39), written in
the gap gauge of \cref{sec:conventions}, is the trace
\begin{equation}
  \mathcal R
  \equiv
  \tfrac14\Tr\!\bigl[\hatD\,(\hatg^R+\hatg^A)\bigr]
  =
  0
  \qquad(\text{ideal BCS, } E>\Delta),
  \label{eq:depairing_coefficient}
\end{equation}
because the propagators satisfy \(\hatg^R+\hatg^A=0\) above the
gap (the opposite advanced-propagator sign would instead give
\(\mathcal R=\Delta N_2\), the anomalous components adding rather than
cancelling); consistently, \cref{eq:fT_kinetic_conservative} contains no
coherent above-gap relaxation. Off the ideal-BCS background
\(\mathcal R\propto\Delta\,\Imag\sinh\theta\) is generically nonzero---under
depairing, proximity, or sub-gap conditions, where \(\theta\) is
complex---and the corresponding relaxation terms enter the transverse
equation through the same channel projection evaluated on the complex
spectra. Likewise a supercurrent couples the two channels through
\(S=2Q\,\Imag\sinh^{2}\theta\) (verified symbolically), which
vanishes on the real-\(\Delta\), zero-\(Q\) background of this section.
These extensions, together with the inelastic content of \(J_T\), are the
standard charge-imbalance relaxation channels; none of them is needed
for, or modified by, the coherent terms of
\cref{eq:fT_kinetic_conservative}.

\paragraph{Spatially varying gap.}
For a static real \(\Delta(\br)\) the kernel-channel flux generalizes
\cref{eq:fT_static} to
\begin{equation}
  \hbar\,\DOS\,\partial_t\fT
  =
  \frac{\hbar\DN}{\DOS}\,
  \widetilde\bnabla\!\cdot\!\bigl[\DOS\,\widetilde\bnabla\fT\bigr]
  +\frac{\hbar\DN}{\DOS}\,N_2\,
  \widetilde\bnabla\!\cdot\!\bigl[\fT\,\widetilde\bnabla\vartheta\bigr]
  +\hbar\,J_T ,
  \qquad
  \vartheta=\operatorname{artanh}\frac{\Delta}{E},
  \label{eq:fT_inhom}
\end{equation}
with \(\widetilde\bnabla\vartheta=E\,\widetilde\bnabla\Delta/W^{2}\) the
pairing-angle gradient: alongside the weighted diffusion, a gap gradient
coherently converts charge imbalance, the transverse counterpart of the
statement that the longitudinal channel has \emph{no} DOS-gradient drift.
The \(\tau_3\)-trace of the same flux is
\(\widetilde\bnabla\!\cdot\!(\DOS^2\widetilde\bnabla\fT)\), the divergence
of the dressed current \(\mathbf j_T\), consistent with
\cref{eq:traceT_diffusion}: through \(\tau_3=\DOS\,\hatg^0-N_2\hat m\),
the dressed-divergence grouping familiar from the static literature is
the \(\DOS\)-weighted kernel channel plus slaved \(\hat m\)-channel
content, not an independent third scalar equation. The slaved
image-channel residual
\(\propto\DOS\,(\widetilde\bnabla^{2}\vartheta)\fT\) remains in the
\(\hat m\) channel and fixes the subleading \(\tau_1\) content of
\(\hatg_0^K\) without entering the scalar dynamics at this order.

\paragraph{Scope.} The transverse sector established here comprises the
kinetic equation
\cref{eq:fT_kinetic_conservative,eq:fT_static,eq:fT_inhom}, the current
dressing \cref{eq:traceT_diffusion}, and the trace statement
\cref{eq:depairing_coefficient}. Outside it remain: gauge-potential
(\(\partial_t\phi\), charge-mode) drives---gap-control problems in the
literature generically carry such terms, which must not be conflated with
a real \(\dot\Delta\); the complex-spectrum (depairing/proximity)
relaxation and sub-gap Andreev charge transport, where the kernel
geometry \cref{eq:kernel_channels} itself deforms; and the microscopic
evaluation of \(J_T\), for which near-\(T_c\) forms exist in the
literature (Kopnin~\S 10.4).

\subsubsection{Scalar reduction: \texorpdfstring{\(\fT\to 0\)}{fT to 0} and the
occupation representation}
\label{sec:scalar_form}

For \emph{static} \(\Delta\), the proven longitudinal equation
\cref{eq:static_kinetic_fL} translates to the scalar occupation
\(f=(1-\fL)/2\) cleanly. With \(\partial_t\fL=-2\partial_t f\) and
\(\widetilde\bnabla\fL=-2\widetilde\bnabla f\), and \(\partial_t\DOS=0\),
substituting into \cref{eq:static_kinetic_fL} and dividing through by
\(-2\) gives
\begin{equation}
  \boxed{\,
  \hbar\,\DOS(E,\br)\,\partial_t f
  =
  \hbar\DN\,\widetilde\bnabla^{2} f
  + \hbar \Icoll^{\mathrm{inel}}[f,n]
  \quad(\text{static }\Delta,\;E>\Delta(\br)),
  \,}
  \label{eq:final_scalar}
\end{equation}
where \(\Icoll^{\mathrm{inel}}=-(1/2)J_1=\DOS\,I_{\mathrm{occ}}\) is the
inelastic collision integral in the occupation variable, still carrying
the external spectral weight (\cref{eq:J1_occ_bridge}); dividing the
equation through by \(\DOS\) yields the bare occupation form
\(\partial_t f=(\DN/\DOS)\widetilde\bnabla^2 f+I_{\mathrm{occ}}\). This
is the scalar dirty-limit
Boltzmann-like kinetic equation for static gap; below the local gap edge
the flux is blocked (\(\mathcal D_L=0\)).

The elastic impurity collision integral does not appear on the right-hand
side: it has been absorbed into the diffusion coefficient \(\DN\) via
the matrix angular reduction in \cref{sec:matrix_usadel}. The only
explicit collision integral is inelastic.

\paragraph{Time-dependent extension.} For time-dependent
\(\Delta(t)\), the conservative form \cref{eq:full_kinetic_conservative}
that the plain-trace projection produces already contains the
\(\partial_E[(\Delta/E)\dot\Delta\,\DOS\fL]\) energy-space flux, so the
fixed-\(E\) trace equation translates directly to the target equation
\cref{eq:target}: with the DOS-continuity identity it reads
\(\DOS\,(\partial_t f+(\Delta/E)\dot\Delta\,\partial_E f)
=\DN\,\widetilde\bnabla^{2} f
+\Icoll^{\mathrm{inel}}\).
The coordinate lift \cref{eq:coordinate_lifted_A1_advective}, the
moving-bin weak projection of
\cref{sec:moving_bin_projection}, and the branch-projector construction of
\cref{sec:branch_projector_lift} reproduce the same equation kinematically.
Whether a representation-level matrix-Usadel reformulation---with the moving
spectral control volume present before the fixed-\(E\) scalar projection is
taken---supplies anything the fixed-\(E\) route does not is settled in
\cref{sec:resolution}: it is a representation choice (rigor/manifestness), not a
source of new terms.

\paragraph{Operator identification.} \cref{eq:final_scalar} has
the form
\(\DOS\,\partial_t f=\DN\,\widetilde\bnabla^{2} f
+\Icoll+\cdots\), or equivalently
\(\partial_t f=L_{D,\mathrm{U}}^{(1,0)}[f]+\DOS^{-1}\Icoll+\cdots\)
with
\begin{equation}
  L_{D,\mathrm{U}}^{(1,0)}[f]
  =
  \frac{\DN}{\DOS(E,\br)}\,
  \widetilde\bnabla\!\cdot\!
  \bigl[\mathcal D_L(E,\br)\,\widetilde\bnabla f\bigr],
  \label{eq:LD_p1q0}
\end{equation}
the \((p,q)=(1,0)\) member of the main text's \(L_{p,q}\) family
(main-text~\cref{M-eq:Lpq_family}), with \(\mathcal D_L=1\) above the local gap edge.
For uniform gap, this collapses to
\(L_{D,\mathrm{U}}^{(1,0)}[f]=(\DN/\DOS(E))\nabla^{2}_{\!\br}f\), an
energy-dependent rate that vanishes at the gap edge as
\(\DOS(E)\to\infty\) and recovers the normal-state value
\(\DN\) far above the gap---the same spectral suppression as the clean
Bardeen--Rickayzen--Tewordt closure. This is the dirty-limit answer.

\paragraph{Ingredients.} The conclusion that the dirty-limit
matrix-Usadel route gives operator A1 \((p,q)=(1,0)\), with conserved
density \(\DOS\fL\) and undressed flux, uses (i) the
spatial-flux derivation of \cref{sec:fLfT} (which fixed \(q=0\) from
the BCS bilinear identities and the matrix-current decomposition), (ii)
the trace identity \cref{eq:trace_identities}
\(\Tr[\tau_3\hatg_0^K]=4\DOS\fL\) (which fixed the conserved spectral
density and hence \(p=1\)), and (iii) the matrix Moyal commutator
\cref{eq:moyal_comm} (which produced the \(\partial_t\) term and the
spectral-flow energy flux on the
LHS). None of these is a free choice; the \((p,q)=(1,0)\) structure is
the unique result of the standard matrix-Usadel reduction within the
plain-trace projection carried out here, for static and uniform
time-dependent gaps alike.

The diagnostic \((p,q)=(2,2)\) operator (main text
main-text~\cref{M-eq:usadel_LD_main}) is \emph{not} generated by this
projection: it would require the conserved density \(\DOS^{2}\fL\) on
the LHS, which the trace identity \cref{eq:trace_identities} does not
support. As the main text notes, A2 is a useful symmetric test
operator, but it is not what the matrix Usadel produces; A1 is.

\subsection{Alternative: branch-Boltzmann route}
\label{sec:branch_alternative}

\subsubsection{Project first, angular-average second}

The scalar route reverses the order of operations used above. Instead of
first impurity-averaging the matrix Green's function into a Usadel object and
only then tracing onto \(\fL\), one first projects the microscopic dynamics onto
the positive-energy Bogoliubov branch. The branch occupation
\(f(E,\hat\bp,\br,t)\) then obeys the semiclassical kinetic equation
\begin{equation}
  \partial_t f
  + \frac{\Delta}{E}\dot\Delta\,\partial_E f
  + \vg(E)\,\hat\bp\cdot\bnabla f
  =
  I_{\mathrm{imp}}[f]
  + I_{\mathrm{inel}}[f],
  \label{eq:branch_boltzmann}
\end{equation}
where
\begin{equation}
  \vg(E)=\vF\sqrt{1-\frac{\Delta^2}{E^2}}
  = \frac{\vF}{\DOS(E)}
  \label{eq:branch_velocity}
\end{equation}
on the BCS branch. The spectral-flow term in
\cref{eq:branch_boltzmann} is elementary branch kinematics: at fixed normal
dispersion \(\xi\), \(E(\xi,\Delta)=\sqrt{\xi^2+\Delta^2}\) gives
\(\dot E=(\Delta/E)\dot\Delta\). This is the same scalar term that the
matrix-Usadel route produces through its fixed-\(E\) projection
(\cref{sec:scalar_reduction}); here it
appears before angular averaging because the branch projection has already been
made. The elastic impurity term is explicit at this stage:
\begin{equation}
  I_{\mathrm{imp}}[f](\hat\bp)
  =
  \avg{
  W_E(\hat\bp,\hat\bp')\,
  \bigl[f(\hat\bp')-f(\hat\bp)\bigr]
  }_{\hat\bp'}.
  \label{eq:branch_imp_collision}
\end{equation}
This is a scalar Boltzmann equation for the branch occupation, not yet a
diffusion equation.

Expand the branch distribution on a fixed positive-energy shell,
\begin{equation}
  f(E,\hat\bp,\br,t)
  =
  f_0(E,\br,t)
  +
  \hat\bp\cdot\mathbf f_1(E,\br,t)+\cdots .
  \label{eq:branch_harmonic}
\end{equation}
For an isotropic elastic kernel, the first angular harmonic relaxes with the
transport rate
\begin{equation}
  \frac{1}{\tautr(E)}
  =
  \avg{
  W_E(\hat\bp,\hat\bp')\,
  \bigl(1-\hat\bp\cdot\hat\bp'\bigr)
  }_{\hat\bp'},
  \label{eq:branch_transport_rate}
\end{equation}
and the first moment of \cref{eq:branch_boltzmann} gives, to leading order in
gradients,
\begin{equation}
  \mathbf f_1
  =
  -\tautr(E)\,\vg(E)\,\bnabla f_0.
  \label{eq:branch_f1_slaving}
\end{equation}
The angular average of the elastic collision integral vanishes, because elastic
potential scattering conserves quasiparticle number on the fixed energy shell.
The zeroth moment also fixes the operator ordering: the streaming term of
\cref{eq:branch_boltzmann} is advective---the derivative acts on \(f\)
alone---so \(\vg\), although it varies in space and time through
\(\Delta(\br,t)\), passes through the angular average as a local scalar,
and the divergence acts only on the slaved harmonic. Substituting
\cref{eq:branch_f1_slaving} into the zeroth moment gives
\begin{equation}
  \partial_t f_0
  + \frac{\Delta}{E}\dot\Delta\,\partial_E f_0
  =
  I_{\mathrm{inel}}[f_0]
  +
  \frac{\vg(E)}{d}\,
  \bnabla\cdot\!\left[
    \vg(E)\,\tautr(E)\,\bnabla f_0
  \right],
  \label{eq:branch_diffusion}
\end{equation}
equivalently, with \(\DOS(E)\vg(E)=\vF\),
\begin{equation}
  \partial_t f_0
  + \frac{\Delta}{E}\dot\Delta\,\partial_E f_0
  =
  I_{\mathrm{inel}}[f_0]
  +
  \frac{1}{\DOS}\,
  \bnabla\cdot\!\left[
    \DOS\,D_{\mathrm B}(E)\,\bnabla f_0
  \right],
  \qquad
  D_{\mathrm B}(E)
  =
  \frac{1}{d}\,\vg^2(E)\,\tautr(E).
  \label{eq:branch_diffusion_coefficient}
\end{equation}
The rate coefficient \(D_{\mathrm B}\) is undetermined until the projected
impurity kernel fixes \(\tautr(E)\); the placement of the spectral weight
\(\DOS\) outside the divergence is not---it is the moment algebra itself.
Writing \(\bnabla\cdot[D_{\mathrm B}\bnabla f_0]\), with the whole
coefficient inside the divergence, discards this weight and corresponds to
treating \(f\) rather than \(\DOS f\) as the conserved shell density.

\subsubsection{Projected impurity kernel and BRT cancellation}

For nonmagnetic scalar disorder the Bogoliubov matrix element between
positive-energy branch states contains the charge-like coherence factor
\begin{equation}
  \mathcal C_{\bp\bp'}
  =
  \bigl(u_{\bp}u_{\bp'}-v_{\bp}v_{\bp'}\bigr)^2 .
  \label{eq:branch_coherence_factor}
\end{equation}
On a bulk BCS energy shell, and for scattering between states with the same
radial branch sign,
\begin{equation}
  u_{\bp}u_{\bp'}-v_{\bp}v_{\bp'}
  \simeq
  u_E^2-v_E^2
  =
  \frac{\xi_E}{E}
  =
  \sqrt{1-\frac{\Delta^2}{E^2}}
  =
  \frac{1}{\DOS(E)} ,
  \label{eq:branch_coherence_on_shell}
\end{equation}
where \(\xi_E=\sqrt{E^2-\Delta^2}\). For the opposite radial branch sign
\(\xi\to-\xi\), \(u\) and \(v\) are interchanged and the same case-I matrix
element vanishes, \(uv-vu=0\); the same-branch channel is therefore the
nonzero bulk contribution. The final-state BCS density of states contributes
one factor of \(\DOS(E)\), while the squared coherence factor contributes
\(\DOS^{-2}(E)\). Relative to the normal-state transport rate \(\tau_0^{-1}\),
the branch elastic transport rate therefore scales as
\begin{equation}
  \frac{1}{\tautr(E)}
  =
  \frac{1}{\tau_0\,\DOS(E)},
  \qquad
  \tautr(E)=\tau_0\,\DOS(E),
  \label{eq:branch_brt_time}
\end{equation}
up to the same angular transport factor that defines \(\tau_0\). Since
\(\vg(E)=\vF/\DOS(E)\), the mean free path is
\begin{equation}
  \elltr(E)=\vg(E)\tautr(E)=\vF\tau_0,
  \label{eq:branch_constant_ell}
\end{equation}
and the branch diffusion coefficient becomes
\begin{equation}
  D_{\mathrm{BRT}}(E)
  =
  \frac{\vF^2\tau_0}{d}\,\frac{1}{\DOS(E)}
  =
  \frac{\DN}{\DOS(E)} .
  \label{eq:branch_brt_diffusion}
\end{equation}
This is the Bardeen--Rickayzen--Tewordt clean-elastic cancellation in the
notation of the main text: the density-of-states enhancement in phase space
is cancelled by one power from the impurity coherence factor, and the remaining
velocity factor leaves a constant normal mean free path. In the ordered
reduction \cref{eq:branch_diffusion_coefficient} this closure gives
\(\DOS\,D_{\mathrm{BRT}}=\DN\), so the branch route lands on the
dirty-limit operator A1 exactly.

If instead one ignores the projected impurity coherence factor and simply sets
\(\tautr(E)=\tau_0\), then
\begin{equation}
  D_{\tau=\tau_0}(E)
  =
  \frac{\vF^2\tau_0}{d}\,\frac{1}{\DOS^2(E)}
  =
  \frac{\DN}{\DOS^2(E)} .
  \label{eq:branch_constant_tau_diffusion}
\end{equation}
This constant-\(\tau\) branch closure is a useful diagnostic but is not the
nonmagnetic BRT elastic result.

\subsubsection{Operator labels and the coincidence with the Usadel route}

The ordered branch equation \cref{eq:branch_diffusion_coefficient} carries
the spectral weight \(\DOS\) outside the divergence, so in the main text's
\(L_{p,q}\) notation it has \(p=1\), with \(q\) fixed by the closure. The
BRT kernel gives \(\DOS\,D_{\mathrm{BRT}}=\DN\),
\begin{equation}
  \DOS\,\partial_t f_0
  =
  \bnabla\cdot\!\left[
    \DN\,\bnabla f_0
  \right]
  + \cdots ,
  \qquad
  (p,q)=(1,0),
  \label{eq:branch_brt_operator}
\end{equation}
identical, above the local gap edge, to the dirty matrix-Usadel result
\begin{equation}
  \DOS\,\partial_t f
  =
  \bnabla\cdot\!\left[
    \DN\,\mathcal D_L\,\bnabla f
  \right]
  + \cdots ,
  \qquad
  (p,q)=(1,0),
  \label{eq:branch_compare_usadel}
\end{equation}
while the constant-\(\tau\) diagnostic gives \((p,q)=(1,-1)\). The
parent-paper rows B \((0,-2)\) and C \((0,-1)\) arise only by moving the
outside spectral weight into the divergence,
\begin{equation}
  \partial_t f_0
  =
  \bnabla\cdot\!\left[
    D_{\mathrm B}(E)\,\bnabla f_0
  \right]
  + \cdots ,
  \label{eq:branch_constant_tau_operator}
\end{equation}
a Fick ansatz for the occupation itself. These legacy placements share the
uniform-gap rates of their ordered counterparts and differ from them, for
an inhomogeneous gap, by the DOS-gradient drift of
main-text~\cref{M-eq:Lpq_expanded}; they are not the end point of the
branch reduction.

The two routes therefore coincide, although the order of operations has
changed. In the branch route, one projects onto a quasiparticle with
branch velocity \(\vg=\vF/\DOS\) and then uses the coherence-dressed
elastic collision integral to slave the first angular harmonic; in the
Usadel route, the impurity self-energy is used first to slave the matrix
first harmonic,
\(\check{\mathbf g}_1=-\ell\checkg_0\widetilde\bnabla\checkg_0\), and the
branch occupation is extracted only afterward by a Nambu trace. The
operations do not commute termwise---the branch-to-trajectory relabelling
acts on odd angular harmonics, so the branch route's current-carrying
first harmonic is the anisotropic \emph{transverse} amplitude of the
two-mode description in disguise---but the relabelling is exhausted by
this exchange of sectors, and the closed longitudinal end points agree:
main-text~\cref{M-eq:projection_average_commutator} marks a boundary
between representations, not between diffusion laws.

Consequently a literal ``dirty scalar Boltzmann equation with
\(I_{\mathrm{imp}}\) left on the right-hand side'' is not an independent
answer. If the elastic impurity term remains explicit, the equation is still
pre-diffusive; solving its first angular moment removes the elastic term from
the isotropic equation and produces \cref{eq:branch_diffusion}. If the system is
already in the dirty Usadel limit, the same elastic scattering has already been
resummed into the matrix angular reduction and into the normal-state scale
\(\DN=\vF^2\tau/d\). Adding a second scalar elastic collision integral would
double-count the impurity physics. Both orderings of the pair of
reductions---project to the branch first and then angular-average, or
angular-average the matrix equation first and then trace/project---give
the static dirty-limit A1 operator derived in \cref{eq:LD_p1q0}, with the
branch-Boltzmann rate coefficients
\cref{eq:branch_constant_tau_diffusion,eq:branch_brt_diffusion} entering
through the ordered form \cref{eq:branch_diffusion_coefficient}.

\subsection{Scalar boundary conditions from the Kupriyanov--Lukichev relation}
\label{sec:boundary_conditions}

The bulk reduction fixed the undressed longitudinal flux
\(\mathbf j_L=-\DN\,\mathcal D_L\bnabla\fL\) and the operator A1,
\cref{eq:LD_p1q0}. Devices are boundary-value problems, so the scalar model also
needs the matching condition at an interface. It follows by projecting the
Kupriyanov--Lukichev (K-L) matrix-current relation~\cite{kupriyanov1988} onto the
same longitudinal occupation used in the bulk.

\paragraph{Matrix relation.} For two diffusive electrodes \(1,2\) joined by a
low-transparency barrier of specific resistance \(R_bA\),
\begin{equation}
  \sigma_i\,\check g_i\,\hat n\!\cdot\!\widetilde\bnabla\check g_i
  =\frac{1}{2R_bA}\,[\check g_1,\check g_2],
  \label{eq:KL_BC}
\end{equation}
so the matrix current through the barrier is the commutator
\(\check I=(2R_bA)^{-1}[\check g_1,\check g_2]\), continuous across the interface.
Take BCS spectral functions on each side,
\(\hatg_i^{R}=\DOS^{(i)}\tau_3+\ii N_2^{(i)}\tau_2\),
\(\hatg_i^{A}=-\DOS^{(i)}\tau_3-\ii N_2^{(i)}\tau_2\), with
\(\DOS^{(i)}=E/\Omega_i\), \(N_2^{(i)}=\Delta_i/\Omega_i\),
\(\Omega_i=\sqrt{E^2-\Delta_i^2}\) and independent gaps \(\Delta_1,\Delta_2\), and
distributions \(\hat h_i=\fL^{(i)}\openone+\fT^{(i)}\tau_3\), so
\(\hatg_i^{K}=\hatg_i^{R}\hat h_i-\hat h_i\hatg_i^{A}\).

\paragraph{Keldysh projection.} The Keldysh component of the commutator is
\begin{equation}
  [\check g_1,\check g_2]^{K}
  =\hatg_1^{R}\hatg_2^{K}+\hatg_1^{K}\hatg_2^{A}
   -\hatg_2^{R}\hatg_1^{K}-\hatg_2^{K}\hatg_1^{A}.
  \label{eq:KL_keldysh}
\end{equation}
Projecting with the same plain Nambu trace that gives the longitudinal flux in
the bulk (and the \(\tau_3\)-trace for the charge channel),
\begin{align}
  \tfrac14\Tr\!\big[[\check g_1,\check g_2]^{K}\big]
    &=-2\,\big(\DOS^{(1)}\DOS^{(2)}-N_2^{(1)}N_2^{(2)}\big)\,(\fL^{(1)}-\fL^{(2)}),
  \label{eq:scalar_BC_L}\\
  \tfrac14\Tr\!\big[\tau_3[\check g_1,\check g_2]^{K}\big]
    &=-2\,\DOS^{(1)}\DOS^{(2)}\,(\fT^{(1)}-\fT^{(2)}).
  \label{eq:scalar_BC_T}
\end{align}
No \(\fT\) enters the energy channel and no \(\fL\) the charge channel: as in the
bulk for a real gap, the interface does not mix longitudinal and transverse
modes. (Both identities are confirmed by symbolic computer algebra.)

\paragraph{The three scalar objects.} Writing \(G_N=(R_bA)^{-1}\), the energy
(longitudinal) interface current is
\begin{equation}
  \boxed{\;
  j_L=G_N\,\big[\DOS^{(1)}\DOS^{(2)}-N_2^{(1)}N_2^{(2)}\big]\,
  \big(\fL^{(1)}-\fL^{(2)}\big)\;}
  \label{eq:scalar_BC_boxed}
\end{equation}
the coherence-factor weight familiar from the Maki--Griffin
(zero-phase-difference) heat current across a tunnel
junction~\cite{maki1965}. This fixes
the three objects a scalar boundary condition
requires:
\begin{itemize}
  \item \emph{interface weight}
  \(\mathcal W=\DOS^{(1)}\DOS^{(2)}-N_2^{(1)}N_2^{(2)}\) (energy) or
  \(\DOS^{(1)}\DOS^{(2)}\) (charge);
  \item \emph{mobility} \(\mathcal D=G_N=(R_bA)^{-1}\);
  \item \emph{jump variable}: the bare occupation \(\fL\), with the matrix
  current --- hence \(j_L\) --- continuous across the interface, \emph{not} \(\fL\)
  or \(\DOS\fL\). Matching \cref{eq:scalar_BC_boxed} to the bulk flux on each side,
  \(-\DN\,\hat n\!\cdot\!\bnabla\fL\big|_i=j_L\), gives a Robin condition: the
  normal derivative of \(\fL\) is tied to the weighted occupation jump. \(\fL\) is
  discontinuous at finite \(G_N\); the transparent limit \(G_N\to\infty\) forces
  \(\fL^{(1)}=\fL^{(2)}\).
\end{itemize}

\paragraph{Coherence factors and limits.} The energy weight
\(\DOS^{(1)}\DOS^{(2)}-N_2^{(1)}N_2^{(2)}
=(E^2-\Delta_1\Delta_2)/\Omega_1\Omega_2\) carries the anomalous
\(N_2^{(1)}N_2^{(2)}\) that suppresses the DOS product; the
charge weight \(\DOS^{(1)}\DOS^{(2)}=E^2/\Omega_1\Omega_2\) is the bare
density-of-states product familiar from SIS quasiparticle tunneling ---
the standard coherence-factor distinction between energy and
charge-imbalance tunneling, with the channel assignments fixed by
\cref{eq:scalar_BC_L,eq:scalar_BC_T}.
For equal gaps \(\mathcal W_L=\DOS^2-N_2^2=1\) (regular: the energy channel
matches the undressed bulk \(\mathcal D_L=1\)) and \(\mathcal W_T=\DOS^2\)
(matching the bulk \(\mathcal D_T=\DOS^2\)); in the normal-contact
limit \(\Delta_2\to0\) both weights reduce to the superconductor density of
states, \(\mathcal W_L=\mathcal W_T=\DOS^{(1)}\); and for
\(E>\max(\Delta_1,\Delta_2)\) both weights are positive
(\(E^2>\Delta_1\Delta_2\)), so the current flows
down the occupation difference.

\paragraph{Gap-edge caveat.} Per energy, both weights inherit an integrable
\((E-\Delta_{\max})^{-1/2}\) BCS edge singularity at the larger gap edge
for \emph{unequal} gaps. At \emph{matched} gaps the two channels separate
sharply: the energy weight is regular, \(\mathcal W_L=1\), while the
charge weight \(\mathcal W_T=\DOS^2\sim(E-\Delta)^{-1}\) carries the
matched-gap (SIS) logarithmic singularity and must be regularized by gap
broadening, self-consistency, or a vanishing imbalance jump at the edge.

This is the interface counterpart of the bulk operator A1: the same spectral
ingredients \(\DOS,N_2\) reappear as interface weights instead of flux dressings.

\subsection{Beyond the real gap: supercurrent-induced longitudinal--transverse
coupling}
\label{app:supercurrent}

The longitudinal--transverse decoupling established above---for the bulk current
\(\tfrac14\Tr[\check J^K]=\bnabla\fL\), which carries no \(\fT\)---is
special to a real gap. A supercurrent reopens the coupling---on the
ideal BCS background, only weakly. With
a phase gradient the gauge-covariant gradient is
\(\widetilde\bnabla\hatg=\bnabla\hatg-\ii Q[\tau_3,\hatg]\), \(Q\) the
gauge-invariant superfluid momentum. For the \emph{frozen} BCS spectral functions
the supercurrent still does \emph{not} mix the channels: with the corrected
propagators the \(\fT\)-sector Keldysh symbol is
\(\hatg^K|_{\fT}=2\DOS\fT\,\openone\), so
\([\tau_3,\hatg^K|_{\fT}]=0\) identically, and the remaining \(Q\)-term
\(\hatg^K|_{\fT}[\tau_3,\hatg^A]=-4\DOS N_2\fT\,\tau_1\) is traceless in both
scalar channels. The longitudinal
current remains the undressed \(\bnabla\fL\) at \emph{all} \(Q\), and the
transverse current \(\DOS^2\bnabla\fT\) (verified symbolically: all
\(Q\)-dependent terms cancel in both projections).

The coupling comes instead from the current-carrying \emph{depaired} spectrum.
The supercurrent enters the retarded Usadel equation as a depairing energy
\(\Gamma=\tfrac12\hbar D q^2=2\hbar D Q^2\)~\cite{belzig1999}, so that with
\(\hatg^R=\tau_3\cosh\theta+\ii\tau_2\sinh\theta\) the spectral angle solves
\begin{equation}
  E\sinh\theta-\Delta\cosh\theta+\ii\Gamma\sinh\theta\cosh\theta=0,
  \qquad \theta\to\theta_{\mathrm{BCS}}\ (Q\to0),
  \label{eq:depairing}
\end{equation}
and becomes complex. The longitudinal--transverse coupling is the spectral
supercurrent
\begin{equation}
  S(E)=\tfrac14\Tr\big[\tau_3(\hatg^R\widetilde\bnabla\hatg^R
       -\hatg^A\widetilde\bnabla\hatg^A)\big]=2Q\,\Imag\,\sinh^2\theta,
  \label{eq:spectral_supercurrent}
\end{equation}
which vanishes identically for the real (BCS) angle. To first order in \(\Gamma\),
with \(\theta=\theta_0-\ii\Gamma\sinh\theta_0\cosh\theta_0/W\) and
\(W=\sqrt{E^2-\Delta^2}\),
\begin{equation}
  \Imag\,\sinh^2\theta=-\frac{2\Gamma E^2\Delta^2}{W^5},
  \qquad
  \boxed{\,S(E;Q)=-\frac{4Q\Gamma E^2\Delta^2}{W^5}
       =-\frac{8\hbar D\,Q^3 E^2\Delta^2}{W^5}\,}+O(Q^5).
  \label{eq:S_result}
\end{equation}
It enters the coupled dirty-limit kinetic equations~\cite{belzig1999} as
\begin{equation}
  \mathbf j_L=-\DN\big(D_L\bnabla\fL+S\,\fT\big),\qquad
  \mathbf j_T=-\DN\big(D_T\bnabla\fT+S\,\fL\big),
  \label{eq:coupled_LT}
\end{equation}
with \(D_L=\tfrac14\Tr(1-\hatg^R\hatg^A)\) and
\(D_T=\tfrac14\Tr(1-\hatg^R\tau_3\hatg^A\tau_3)\) (reducing to \(1\) and \(\DOS^2\)
in the BCS limit), so a phase gradient mixes the energy (\(\fL\)) and charge
(\(\fT\)) modes. The mixing is third order in the superfluid momentum---weak at
small currents, growing toward the critical current---and vanishes as \(Q\to0\),
consistent with the real-gap decoupling. The \(O(Q^3)\) onset is specific
to the ideal bulk BCS background, where the spectral angle is real at
\(Q=0\) and an imaginary part is generated only by the \(O(Q^2)\)
depairing; on a background whose zero-current spectrum is already
complex---proximity, depairing of other origin, or sub-gap
energies---\(S(E)=2Q\,\Imag\sinh^2\theta_0(E)+O(Q^3)\) is generically
\emph{first} order in \(Q\). (Cancellation and \cref{eq:S_result}
verified symbolically and numerically by computer algebra.)


\subsection{Beyond the real gap II: non-BCS and proximity spectra}
\label{app:proximity}

The supercurrent coupling of \cref{app:supercurrent} required a phase gradient. A
complementary generalisation keeps the order parameter \emph{real} but lets the
spectrum leave the clean BCS form: near a normal-metal contact, inside a proximity
minigap, or under depairing, the spectral angle
\(\theta(E,\br)=\theta'+\ii\theta''\) becomes complex and varies in space while
carrying no current. This is the follow-up calculation flagged below
\cref{eq:flux_dressings_BCS}. The bulk operator A1 survives it intact: the
energy-channel coefficient is promoted from the on/off indicator
\(\mathcal D_L\) to a smooth, manifestly real function of the complex
angle, the energy and charge channels stay decoupled, and the only new ingredient is
a genuine spectral-gradient drift in the energy channel.

Off the real gap the retarded--advanced symmetry must be stated explicitly. In the
\(\ii\tau_2\) gauge (\(\hatg^R=\tau_3\cosh\theta+\ii\tau_2\sinh\theta\)) it is the
\(\tau_3\)-sandwiched Hermitian conjugate of
main-text~\cref{M-eq:bcs_advanced_matrix},
\begin{equation}
  \hatg^A=-\tau_3(\hatg^R)^{\dagger}\tau_3
         =-\cosh\theta^{\ast}\,\tau_3-\ii\sinh\theta^{\ast}\,\tau_2,
  \label{eq:advanced_proximity}
\end{equation}
fixed by the branch structure \(W_A=W_R^{\ast}\) of
\(W_R=\sqrt{(E+\ii0)^{2}-\Delta^{2}}\) together with the large-energy limit
\(\hatg^A\to-\tau_3\). For real \(\theta\) it reduces to the
\(\hatg^A=-\DOS\,\tau_3-\ii N_2\,\tau_2=-\hatg^R\) used throughout
\cref{app:derivation}. Two other continuations must be avoided: the
algebraic shortcut \(-\tau_3\hatg^R\tau_3\) carries \(\theta\) instead of
\(\theta^{\ast}\) and returns complex coefficients, while the plain
conjugate \(-(\hatg^R)^{\dagger}\) swaps the two channel weights below;
the two coincide on the real-\(\theta\) axis.

With \cref{eq:advanced_proximity} the two spectral diffusion coefficients of
\cref{eq:coupled_LT} evaluate, for \emph{any} complex \(\theta\), to
\begin{equation}
  D_L=\tfrac14\Tr\!\big[\openone-\hatg^R\hatg^A\big]=\cos^{2}\theta'',
  \qquad
  D_T=\tfrac14\Tr\!\big[\openone-\hatg^R\tau_3\hatg^A\tau_3\big]=\cosh^{2}\theta',
  \label{eq:DLDT_proximity}
\end{equation}
both manifestly real, with \(\theta'=\Real\theta\) and \(\theta''=\Imag\theta\).
Since the density of states is
\(\DOS=\Real\cosh\theta=\cosh\theta'\cos\theta''\), they obey the exact identity
\begin{equation}
  D_L\,D_T=\DOS^{2}.
  \label{eq:DLDT_identity}
\end{equation}
The longitudinal (energy) channel carries \(D_L=\cos^{2}\theta''\), reducing to
\(1\) on the real gap (\(\theta''=0\)) and to \(0\) deep in the sub-gap
regime (\(\theta''\to\pi/2\): energy transport blocked, in agreement with
Belzig et~al.~\cite{belzig1999}); the transverse (charge) channel
carries \(D_T=\cosh^{2}\theta'\to\DOS^{2}\) on the real gap and remains
nonzero sub-gap (Andreev charge penetration). The product
\cref{eq:DLDT_identity} ties both
dressings to the DOS for an arbitrary spectrum, and \cref{eq:DLDT_proximity} reduces
to \cref{eq:flux_dressings_BCS} in the BCS limit.

Reducing the Keldysh spatial current \(\check J^{K}=(\checkg\,\bnabla\checkg)^{K}\)
for a complex, spatially varying \(\theta(E,\br)\) with
\cref{eq:advanced_proximity} gives
\begin{equation}
  \tfrac14\Tr[\check J^{K}]=\cos^{2}\!\theta''(E,\br)\,\bnabla\fL,
  \qquad
  \tfrac14\Tr[\tau_3\check J^{K}]=\cosh^{2}\!\theta'(E,\br)\,\bnabla\fT,
  \label{eq:proximity_currents}
\end{equation}
with \emph{no} cross term and \emph{no} \(\bnabla\theta\) source: a real proximity
spectrum dresses the two channels unequally but leaves them decoupled, exactly as on
the real gap. The spectral-gradient terms cancel in the projected current, which
completes the spatially varying reduction flagged below \cref{eq:flux_dressings_BCS}.
The energy-resolved operator A1 therefore generalises to
\begin{equation}
  \partial_t\!\big(\DOS\,\fL\big)
    =\bnabla\!\cdot\!\Big[\DN\,\cos^{2}\!\theta''(E,\br)\,\bnabla\fL\Big],
  \label{eq:A1_proximity}
\end{equation}
with the conserved density still \(\DOS\,\fL\). The one qualitatively new feature
relative to the uniform gap is that \(\bnabla\theta\neq0\) is now physical rather
than a relabelling: the dressing varies in space, so \cref{eq:A1_proximity} carries a
real drift \(\DN\,[\bnabla\cos^{2}\theta''(E,\br)]\!\cdot\!\bnabla\fL\) absent for a
real-gap spectrum (where \(\theta''\) jumps between \(0\) and its sub-gap
value only at the local gap edge). Longitudinal--transverse mixing reappears only through a
supercurrent (\cref{app:supercurrent}), never through the proximity gradient alone.

The reductions \cref{eq:DLDT_proximity}--\cref{eq:proximity_currents}, the identity
\cref{eq:DLDT_identity}, and the absence of channel mixing and of a
spectral-gradient source were verified symbolically and numerically by computer
algebra for generic complex \(\theta(E,\br)\).


\subsection{Beyond the adiabatic limit: nonadiabatic spectral dynamics}
\label{app:nonadiabatic}

The time-dependent spectral flow of main-text~\cref{M-sec:tdep_spectral_flow} is the
\emph{adiabatic} reduction: it is exact to
first Wigner (Moyal) order, the instantaneous BCS propagator solving
\([L_0,\hatg^R]_\star=0\) at \(O(\hbar^0)\) and \(O(\hbar^1)\). The last
beyond-assumption thread is the leading nonadiabatic correction, obtained from
the second-order Moyal expansion on
the \((E,t)\) phase space.

\paragraph{Spectrum.} Carrying \([L_0,\hatg^R_0]_\star\) for the instantaneous BCS
propagator \(\hatg^R_0=\DOS\tau_3+\ii N_2\tau_2\) to second order, the \(O(\hbar^0)\)
and \(O(\hbar^1)\) terms vanish and the leading residual is purely off-diagonal,
\begin{equation}
  [L_0,\hatg^R_0]_\star
  =\frac{3\hbar^2 E\,\Delta^2\,\ddot\Delta}{4\,W^5}\,\tau_1+O(\hbar^3),
  \qquad W=\sqrt{E^2-\Delta^2}.
  \label{eq:nonad_residual}
\end{equation}
The instantaneous spectral functions are therefore exact through \(O(\hbar)\); the
nonadiabatic correction to \(\hatg^R\) enters only at \(O(\hbar^2\ddot\Delta)\) and in
the \(\tau_1\) (branch) sector, so the density of states and the spectral
coefficients \(\mathcal D_L,\mathcal D_T\) are robust against a
time-dependent gap to this order.

\paragraph{Distribution.} Form \(\hatg^K=\hatg^R_0\star\hat h-\hat h\star\hatg^A_0\)
with \(\hat h=\fL\openone+\fT\tau_3\) and
\(\hatg^A_0=-\tau_3(\hatg^R_0)^\dagger\tau_3=-\hatg^R_0\). The time-dependent kinetic
operator \(\ii[L_0,\hatg^K]_\star\) keeps the energy and charge modes \emph{decoupled}:
through \(O(\hbar^2)\) the energy projection \(\tfrac14\Tr[\,\cdot\,]\) carries no
\(\fT\) and the charge projection \(\tfrac14\Tr[\tau_3\,\cdot\,]\) no \(\fL\),
with the energy projection reproducing the conservative spectral flow of
\cref{eq:expected_conservative} exactly at \(O(\hbar)\) and receiving no
\(O(\hbar^2)\) correction. A
time-dependent gap thus opens \emph{no} direct longitudinal--transverse coupling, in
contrast to the supercurrent of \cref{app:supercurrent}. The off-diagonal
(branch-coherence) projection vanishes at first order,
\begin{equation}
  \tfrac14\Tr\big[\tau_1\,\ii[L_0,\hatg^K]_\star\big]\Big|_{O(\hbar)}
  =0
  \label{eq:nonad_coherence}
\end{equation}
(a spurious \(O(\hbar\dot\Delta)\) drive appears here only for the opposite
advanced-propagator sign). The leading branch-coherence generation appears
at \(O(\hbar^2)\), sourced by the energy mode alone with coefficients
proportional to \(\ddot\Delta\) and \(\dot\Delta^2\)---the same order as the
retarded spectral residual \cref{eq:nonad_residual}. Being off-diagonal it
does not enter the conserved scalar densities; its relaxation feeds the
collision-integral sector rather than \(\fL\) directly.

\paragraph{Summary.} A fast gap leaves the energy and charge diffusion modes decoupled
at the spectral-flow level and the spectrum adiabatic through \(O(\hbar)\); \emph{all}
leading nonadiabaticity---the branch-coherence drive and the spectral
residual \cref{eq:nonad_residual} alike---enters at
\(O(\hbar^2)\), through \(\ddot\Delta\) and \(\dot\Delta^2\), and vanishes as
\(\dot\Delta,\ddot\Delta\to0\), recovering main-text~\cref{M-eq:usadel_A1_tdep_fL}.
These are pointwise statements away from the ideal gap edge: the
residuals carry inverse powers of \(W=\sqrt{E^2-\Delta^2}\)
[\cref{eq:nonad_residual}], so the adiabatic ordering requires the
spectral angle to change slowly on the local branch-splitting scale,
\(\hbar|\partial_t\theta|/W\sim\hbar E|\dot\Delta|/W^3\ll1\) (with the
analogous \(\ddot\Delta\) condition); sufficiently close to the edge the
expansion must be regularized by lifetime broadening, self-consistency,
or an explicit edge treatment.
\Cref{eq:nonad_residual}, the energy/charge decoupling, and \cref{eq:nonad_coherence}
were verified symbolically and numerically by computer algebra for generic
\(\Delta(t)\).


\section{Branch-covariant verification}\label{app:branch}


This appendix establishes by explicit Pauli-trace algebra that the
fixed-$E$, fixed-gauge Wigner${}+{}$SSLO Keldysh--Usadel reduction generates
the linear-in-$\dot\Delta$ spectral-flow term directly through the local
plain-trace projection, in the canonical conservative form, and verifies
the companion spatial statements. \Cref{sec:exact} states the exact
identity connecting the conservative and advective (moving-shell) forms;
\cref{sec:plain,sec:spec} compute the projection and the spectral
identities behind it; \cref{sec:fT} bounds the transverse feedback at
$O(\dot\Delta^{2})$; \cref{sec:verdict} collects the interpretation; and
\cref{sec:resolution} verifies, for a spatially varying gap, that the matrix
reduction reproduces the dirty-limit spatial current cleanly, with no spurious
drift and no transverse mixing. The appendix also records the convention
check that fixes the advanced propagator: the opposite sign choice
$-\hatg^{R\dagger}$ fails the advanced spectral equation and removes the
equilibrium anomalous Keldysh weight (\cref{sec:spec}), and with it the
spectral-flow drive. The moving-shell identity supplies the kinematic
reading of the fixed-$E$ result throughout.

\subsection{Setup and projection convention}
\label{sec:setup}

Gap gauge (as in the main text): $\hat L_0=E\tau_3+\ii\Delta(t)\tau_2$,
\begin{equation}
  \hatg^{R}=\DOS\tau_3+\ii N_2\tau_2,\quad
  \hatg^{A}=-\tau_3\hatg^{R\dagger}\tau_3=-\DOS\tau_3-\ii N_2\tau_2=-\hatg^{R},\quad
  \DOS^2-N_2^2=1,\quad N_2=\tfrac{\Delta}{E}\DOS,
  \label{eq:gRA}
\end{equation}
all real above the gap. Distribution $\hat h=\fL\openone+\fT\tau_3$, so
\begin{equation}
  \hatg_0^K=\hatg^{R}\hat h-\hat h\hatg^{A}
  =(2\DOS\tau_3+2\ii N_2\tau_2)\,\fL+2\DOS\,\fT\,\openone
  =2\fL\,\hatg^{R}+2\DOS\,\fT\,\openone .
  \label{eq:gK}
\end{equation}
We use the Moyal commutator
$[A,B]_\star=[A,B]+\tfrac{\ii\hbar}{2}(\partial_EA\,\partial_tB-\partial_tA\,\partial_EB
-\partial_EB\,\partial_tA+\partial_tB\,\partial_EA)+O(\hbar^2)$.

\paragraph{Projection convention.} The \emph{longitudinal} ($\fL$) kinetic equation is the
\textbf{plain} Nambu trace $\Tr[\,\cdot\,]$ of the $K$-block, not the
$\tau_3$-trace. The $\tau_3$ that produces $\partial_t(\DOS\fL)$ comes from the
Moyal factor $\partial_E\hat L_0=\tau_3$ in $\ii[E\tau_3,\hatg_0^K]_\star
=\ii E[\tau_3,\hatg_0^K]-\tfrac{\hbar}{2}\{\tau_3,\partial_t\hatg_0^K\}$, whose
plain trace is $-\hbar\Tr[\tau_3\partial_t\hatg_0^K]=-4\hbar\,\partial_t(\DOS\fL)$.
(The base of this projection is the identity, not $P_+$ or $\tau_3$.)

\subsection{The exact target identity}
\label{sec:exact}

The canonical (branch-Boltzmann) form and the matrix conservative form differ by
an exact identity. With $u_E\equiv(\Delta/E)\dot\Delta$ and
$u_E\DOS=\dot\Delta N_2$,
\begin{equation}
  \boxed{\;
  \partial_t(\DOS\fL)+\partial_E(\dot\Delta\,N_2\,\fL)
  =\DOS\bigl[\partial_t+u_E\partial_E\bigr]\fL\;}
  \label{eq:exact}
\end{equation}
using $\partial_t\DOS=-\partial_E(u_E\DOS)=-\dot\Delta\,\partial_EN_2$ (BCS
DOS-continuity). \textbf{The energy-space flux
$\partial_E(\dot\Delta\,N_2\,\fL)$, which lives entirely in the anomalous
spectral weight $N_2$, is the term that converts between the two forms};
\cref{sec:plain} shows the projection produces it.

\subsection{Direct plain-trace projection: the conservative flow, no
\texorpdfstring{$\fT$}{fT} coupling}
\label{sec:plain}

Compute $\ii[\hat L_0,\hatg_0^K]_\star$ directly and take the plain trace.

\emph{Longitudinal piece} ($\hatg_0^{K}\!\supset 2\fL\,\hatg^R$): since the
symbol is proportional to $\hatg^R\propto\hat L_0$, the ordinary commutator
vanishes identically and the first Moyal order is purely scalar,
\begin{equation}
  \ii[\hat L_0,2\fL\,\hatg^{R}]_\star
  =-2\hbar\bigl[\partial_t(\DOS\fL)+\partial_E(\dot\Delta\,N_2\,\fL)\bigr]
  \openone ,
  \label{eq:Lpiece}
\end{equation}
proportional to the identity with \emph{no} off-diagonal residue (the
equality uses DOS-continuity to absorb the
$\fL(\partial_t\DOS+\dot\Delta\partial_EN_2)$ term). Plain trace:
$-4\hbar\,[\partial_t(\DOS\fL)+\partial_E(\dot\Delta N_2\fL)]$ --- the
complete conservative spectral flow, including the energy-space flux of
\cref{eq:exact}. The $\tau_3$ part of the symbol supplies the
$\partial_t(\DOS\fL)$ term through
$\{\tau_3,\partial_t(\cdot)\}$, and the $\tau_2$ part---present because
$\hatg^R-\hatg^A=2\DOS\tau_3+2\ii N_2\tau_2$ retains the anomalous
weight---supplies $\partial_E(\dot\Delta N_2\fL)$ through
$\{\tau_2,\partial_E(\cdot)\}$.

\emph{Transverse piece} ($\hatg_0^{K}\!\supset 2\DOS\fT\,\openone$): the
ordinary commutator with a scalar symbol vanishes, and the first-Moyal
terms land on $\tau_2,\tau_3$ (none on $\openone$ or $\tau_1$):
\begin{equation}
  \Tr\bigl[\ii[\hat L_0,\hatg_0^{K}|_{\fT}]_\star\bigr]=0
  \qquad(\text{leading and }O(\hbar)).
  \label{eq:fTtrace}
\end{equation}
\textbf{Result.} The longitudinal equation is
$\hbar[\partial_t(\DOS\fL)+\partial_E(\dot\Delta N_2\fL)]
=\hbar\DN\bnabla^{2}\fL+\hbar J_1$,
the fixed-$E$ conservative spectral-flow form, with \emph{no} $\fT$ source,
to first Moyal order --- with $\fT$ retained, not just at $\fT=0$. The drive,
$\hbar\dot\Delta\,\partial_E\fL\,\tau_2$, sits in the $\tau_2$
(anomalous) channel and survives the projection: the sandwich weight is
$\hatg^R\tau_2-\tau_2\hatg^A=2\ii N_2\openone\neq0$ (identity I3 of the
symbolic verification; it vanishes for the opposite advanced-propagator
sign, losing the flow).

\subsection{Spectral identities behind the kinematic reading}
\label{sec:spec}

A clean by-product, computed to first Moyal order:
\begin{equation}
  [\hat L_0,\hatg^{R}]_\star=\ii\hbar\bigl(\partial_t\DOS+\dot\Delta\,\partial_EN_2\bigr)\openone=0,
  \label{eq:specR}
\end{equation}
the vanishing being \emph{identical} to BCS DOS-continuity. So the identity that
powers the kinematic (moving-shell) reading is exactly the statement that
the instantaneous BCS
propagator solves the adiabatic retarded Usadel equation. With
$\hatg^A=-\hatg^R$ the same statement holds
on the advanced side,
\begin{equation}
  [\hat L_0,\hatg^{A}]_\star=0 ,
  \label{eq:specA}
\end{equation}
as it must---retarded and advanced propagators solve the same static
spectral equation---so the clean ``sandwich'' factorization
$[\hat L_0,\hatg^K]_\star=\hatg^R\star[\hat L_0,\hat h]_\star-[\hat L_0,\hat h]_\star\star\hatg^A$
\textbf{holds} to this order, with no $-\hat h\star[\hat L_0,\hatg^A]_\star$
remainder. (The opposite sign choice $-\hatg^{R\dagger}$ gives
$[\hat L_0,\hatg^{A}]_\star=4EN_2\tau_1+2\ii\hbar\dot\Delta\,\partial_EN_2\,\openone\neq0$,
failing the advanced spectral equation; the same choice removes the
anomalous weight from
$\hatg^R-\hatg^A$, forcing $\Delta=0$ in
the equilibrium gap equation. These two requirements are what fix
\cref{eq:gRA}.)

\subsection{The transverse channel enters at \texorpdfstring{$O(\dot\Delta^2)$}{O((dDelta/dt)\string^2)}}
\label{sec:fT}

Could the linear-in-$\dot\Delta$ flow in the energy equation instead be
mediated by charge imbalance? No, by order counting that is independent of
the detailed transverse-sector terms (\cref{sec:transverse_projection}). First,
the plain-trace energy projection carries no direct $\fT$ source at the
orders computed: \cref{eq:fTtrace} at first Moyal order, and the symbolic
verification extends the statement through $O(\hbar^2)$. Second, any
indirect path requires $\fT$ to be generated by the gap drive and then fed
back. The transverse equation of \cref{sec:transverse_projection} makes the
first step explicit: the charge mode is a \emph{linear response} to the drive,
$\fT=O(\dot\Delta)$, with the response coefficient---not the
order---set by the relaxation physics; and every coupling
returning $\fT$ to the energy mode at this level itself vanishes at
$\dot\Delta=0$ (the static channels are exactly decoupled,
\cref{eq:fTtrace,eq:I7,eq:I8}). The feedback is therefore
$O(\dot\Delta^{2})$. \textbf{The spectral-flow term is linear in
$\dot\Delta$; charge imbalance can only feed back at second order.}
Physically: $\fT$ is the linear response to the drive, and its back-action
on the energy mode is one more power of the drive. Hence the linear flow is
necessarily a \emph{direct, kinematic} effect on $\fL$---exactly the term
the projection produces in \cref{eq:Lpiece}---not a
charge-imbalance--mediated one.

\subsection{Interpretation}
\label{sec:verdict}

The target identity holds. Combined with \cref{eq:exact} and
\cref{eq:specR}, the reduction is interpreted as follows:

\begin{quote}
The fixed-$E$, fixed-gauge Wigner $+$ SSLO matrix Keldysh--Usadel reduction
produces the canonical conservative spectral-flow law
$\partial_t(\DOS\fL)+\partial_E(\dot\Delta N_2\fL)=\cdots$
\emph{completely and correctly}, through the local plain-trace projection,
because the corrected retarded--advanced difference retains the anomalous
$\tau_2$ sector that carries the $\dot\Delta\,\partial_E$ drive. The
moving-shell (fixed-$\xi$) form is related to it by the exact identity
\cref{eq:exact}: it is the same law re-expressed on quasiparticle states
rather than Wigner-energy bins, and supplies the kinematic interpretation,
not a missing ingredient.
\end{quote}

A moving-$\xi$ or gauge-invariant-energy reformulation of the starting
representation remains available; it makes the spectral flow manifest
before the scalar limit, and \cref{sec:resolution} shows that it yields no
term the fixed-$E$ reduction lacks.

\paragraph{Caveats.} All traces are first Moyal order, real gap, spatially
uniform, bulk BCS, $\fT$ adiabatic. The transverse kinetic equation is
assembled in \cref{sec:transverse_projection}; \cref{sec:fT} itself uses only
order counting. The trace identities
\cref{eq:Lpiece,eq:fTtrace,eq:specR,eq:specA} have been verified independently by
symbolic computer algebra (identities I1--I4 of the verification
scripts; see the data and code availability statement of the main text).

\subsection{Representation independence for a spatially varying gap}
\label{sec:resolution}

With the time-dependent flow produced directly by the fixed-$E$ projection
(\cref{sec:plain}), the remaining representation question is spatial: does a
moving-$\xi$ / branch-covariant frame generate any spatial term---a
$\bnabla\Delta$ drift or cross-channel current---that the fixed-$E$
reduction lacks? It does not: the spatially varying calculation below
produces no such term, by direct Pauli-trace evaluation and, independently,
by symbolic computer algebra.

\paragraph{Spatially varying gap.} A genuine derivation must in particular reproduce
the dirty-limit \emph{spatial} current in the \emph{same} representation for a
spatially varying gap $\Delta(\br)$, with every $\bnabla\Delta$
spectral-gradient piece produced and cancelled within the representation
itself rather than imported from the static result. We specialize to one spatial dimension $x$ for brevity; the conclusion is identical for the full gradient $\bnabla$. Take the static, real $\Delta(x)$, leading-gradient
Keldysh--Usadel current
\begin{equation}
  \check J^{K}=\hatg^{R}\,\partial_x\hatg^{K}+\hatg^{K}\,\partial_x\hatg^{A},
  \qquad
  \hatg^{K}=\hatg^{R}\hat h-\hat h\hatg^{A},\quad
  \hat h=\fL\openone+\fT\tau_3 .
\end{equation}
With $a\equiv\DOS$, $b\equiv N_2$ (so $a^2-b^2=1$) one has
$\hatg^{K}=2a\fL\tau_3+2\ii b\fL\tau_2+2a\fT\openone$, and projecting,
\begin{align}
  \tfrac14\Tr\!\big[\check J^{K}\big]       &= \partial_x\fL,
  \label{eq:I7}\\
  \tfrac14\Tr\!\big[\tau_3\check J^{K}\big] &= \DOS^{2}\,\partial_x\fT .
  \label{eq:I8}
\end{align}
\textbf{Every $\partial_x\Delta$ term cancels} in both projections: no
$\partial_x\Delta\,\fL$ or $\partial_x\Delta\,\fT$ source, no
$\partial_x\Delta\,\partial_E f$ drift, and \emph{no $\fL$--$\fT$ cross-channel
current} at this order (the algebra keeps $\fT$ arbitrary, so a finite $\fT$ does
not reopen the mixing). The longitudinal current is exactly the A1
flux $\mathbf j_L=-\DN\bnabla\fL$, undressed and \emph{local in $E$}, with
\emph{no} DOS-gradient drift: for the energy channel the only spatial effect
of $\bnabla\Delta$ is the motion of the local gap edge inside
$\mathcal D_L(E,\br)$. The $\DOS^2$ dressing, and with it the
$(\partial_x\DOS^2)\partial_x\fT$ drift inside the divergence, belongs to
the transverse channel \cref{eq:I8}.

\paragraph{Interpretation.} The $\bnabla\Delta\,\partial_E f$ term a \emph{naive}
moving-$\xi$ current would generate via
$\bnabla_\xi F=\bnabla_E f+(\Delta/E)\bnabla\Delta\,\partial_E f$ is simply
\emph{absent} from the fixed-$E$ matrix current: it is a fixed-$\xi$ relabeling
artifact, exactly as in the time case (\cref{eq:exact}). The matrix theory
``derives, not staples'' the A1 spatial current, with no spurious drift and no
transverse mixing.

\paragraph{Collisions.} The kinematic identity \cref{eq:exact} fixes only the
left-hand side; the source must transform too. Within the adiabatic
instantaneous-kernel theory it does, as bookkeeping: if $C_E$ is the fixed-$E$
source density then the per-state/moving-shell source is $C_\xi=C_E/\DOS$, and
instantaneous Fermi-golden-rule kernels (instantaneous DOS, coherence factors,
thresholds such as $2\Delta(t)$) realize exactly this Jacobian weighting. A
genuine inequivalence appears only \emph{outside} the approximation
(nonadiabatic collision memory, drive sidebands, branch coherence, or a gap
changing appreciably during a collision event).

\paragraph{Conclusion.} Under the stated assumptions --- real BCS gap
(uniform or spatially varying), adiabatic, charge-balanced, leading dirty-limit
gradient order, no phase gradient/supercurrent --- the representation question is
\textbf{closed}, and the resolution is negative for ``new physics.'' The two
exact identities
\begin{equation}
  \tfrac14\Tr[\check J^{K}]=\partial_x\fL,
  \qquad
  \partial_t(\DOS\fL)+\partial_E(\dot\Delta\,N_2\fL)
  =\DOS\big[\partial_t+(\Delta/E)\dot\Delta\,\partial_E\big]\fL
\end{equation}
fix the undressed spatial A1 current and the conservative/advective
spectral-flow equivalence, while
collisions transform by the same DOS Jacobian. There is no derivation-level room
for the A2 density $\DOS^2 f$ (unsupported by the trace identity) or for a hidden
$\bnabla\Delta$ longitudinal--transverse mixing. The moving-$\xi$ /
branch-covariant reformulation is a \emph{rigor/manifestness} gain, not a source
of new terms.

\paragraph{Beyond-assumption extensions and their status.}
(i)~A complex order parameter / phase gradient / supercurrent reopens
$\fL$--$\fT$ mixing through the spectral-supercurrent terms; its leading
content is derived in \cref{app:supercurrent}
($S=2Q\,\Imag\sinh^2\theta$, entering through the depaired spectrum,
$O(Q^3)$ on the ideal BCS background and $O(Q)$ on already-complex ones).
(ii)~Non-BCS proximity spectra are treated in \cref{app:proximity}
($\mathcal D_L=\cos^2\theta''$, $\mathcal D_T=\cosh^2\theta'$, no
longitudinal--transverse mixing from a real spectral gradient alone);
higher-gradient corrections remain beyond the leading dirty-limit order.
(iii)~Nonadiabatic \emph{spectral} dynamics is bounded at $O(\hbar^2)$ in
\cref{app:nonadiabatic}; what remains genuinely open is the
collision-level physics (memory, drive sidebands, branch-coherence
relaxation, a gap changing appreciably during a collision) together with
self-consistent gap dynamics. The benchmark problems where A1, A2, and the
legacy scalar placements make distinct predictions are realized in
main-text~\cref{M-fig:bench_uniform,M-fig:bench_drift,M-fig:bench_interface}. Every identity used
in this section is verified independently by the symbolic computer-algebra
scripts cited in the main text.

\bibliographystyle{apsrev4-2}
\bibliography{refs}

\end{document}
